% Chapter 10 -- How Predictable Is Volatility?
% Requires: amsmath, amssymb, booktabs, array, tabularx, graphicx, tcolorbox (breakable, skins), hyperref or url.
% Defines the callout box only if the preamble has not.
\providecolor{fenavy}{HTML}{1B2A6B}
\providecolor{feshade}{HTML}{F3F5FA}
\makeatletter
\@ifundefined{callout}{%
  \newtcolorbox{callout}{breakable, enhanced,
    colback=white, colframe=gray!70, boxrule=0.5pt, arc=1mm,
    borderline west={2.5pt}{0pt}{gray!70},
    left=8pt, right=8pt, top=4pt, bottom=4pt}%
}{}
% Shaded box for optional technical material that general readers may skip.
\@ifundefined{technote}{%
  \newtcolorbox{technote}{breakable, enhanced,
    colback=feshade, colframe=fenavy!60, boxrule=0.5pt, arc=1mm,
    left=8pt, right=8pt, top=4pt, bottom=4pt}%
}{}
\makeatother

\chapter{How Predictable Is Volatility?}
\label{ch:07}
\providecommand{\term}[1]{\emph{#1}}

Anyone who follows the news knows that the stock market has calm stretches, when prices barely move for weeks, and stormy ones, when headlines report large swings day after day. This chapter asks whether that everyday impression can be turned into a reliable forecast of how much prices will move tomorrow.

In 1961 Benoit Mandelbrot, a mathematician at IBM, visited the Harvard economist Hendrik Houthakker and saw a graph of cotton price changes on his blackboard. Houthakker had been unable to fit the data to the normal distribution, the familiar bell curve and the standard model of price changes since Bachelier (1900). Mandelbrot found why: extreme changes were far more common than the normal distribution allows. He also noticed that large changes tend to be followed by large changes, of either sign, and small changes by small changes (Mandelbrot 1963). \textbf{A quiet day tends to follow another quiet day, and a violent day another violent one;} the 1987 crash, the autumn of 2008 and March 2020 each arrived as a run of violent days. Yet for two decades the pattern had no model.

The model came from an unexpected direction. Robert Engle, a physicist turned economist, was on sabbatical at the London School of Economics in 1979, studying inflation rather than markets. Friedman (1977) had argued that the uncertainty created by inflation harms the economy, but standard models assumed that this uncertainty never changes. Engle built a model in which today's variance depends on yesterday's \textbf{squared shocks} (Figure~\ref{fig:10-cluster-arch}). His colleague David Hendry named it \term{autoregressive conditional heteroskedasticity}, or \term{ARCH} (Engle 1982), a long name for a simple idea: the variance, or heteroskedasticity, changes over time in a way that depends on recent events, and it later earned Engle the 2003 Nobel Prize.

ARCH needed many lags to capture volatility that stays high for weeks. Engle's doctoral student Tim Bollerslev, who had come to San Diego from Denmark, solved this by adding yesterday's variance to the equation, so that one term carries the memory of all past shocks (Bollerslev 1986). His \textbf{generalized ARCH}, or \term{GARCH}, became the standard model of financial volatility (Figure~\ref{fig:10-cluster-garch}).


\begin{figure}[htbp]
\centering
\scalebox{0.9}{%
\begin{tikzpicture}[font=\small]
\begin{axis}[width=0.78\linewidth, height=0.34\linewidth,
  axis lines=left, axis line style={-stealth},
  xmin=0, xmax=1000, ymin=-9, ymax=9,
  xtick={0,200,...,1000}, ytick={-8,-4,0,4,8},
  ylabel={Return (\%)}, clip=false,
  xlabel={Trading day},
  legend style={draw=none, fill=none, font=\footnotesize, at={(0.5,-0.45)}, anchor=north,
                legend columns=-1, /tikz/every even column/.append style={column sep=0.8em}},
  legend cell align=left,]
  \addplot[fenavy!80, thin, no marks] coordinates {(0,-0.97) (1,0.28) (2,-0.09) (3,-0.77) (4,1.55) (5,0.01) (6,0.31) (7,-0.77) (8,-0.46) (9,-0.53) (10,0.44) (11,-0.70) (12,0.18) (13,0.39) (14,-0.99) (15,-0.25) (16,-0.23) (17,-0.06) (18,0.61) (19,0.66) (20,0.00) (21,-0.11) (22,1.01) (23,0.01) (24,-0.39) (25,0.06) (26,0.11) (27,0.32) (28,-0.33) (29,0.35) (30,-0.43) (31,0.46) (32,-1.15) (33,0.00) (34,-0.82) (35,0.37) (36,1.12) (37,-2.00) (38,0.60) (39,-0.37) (40,-0.06) (41,1.03) (42,0.23) (43,-1.16) (44,-0.35) (45,-1.63) (46,0.87) (47,-0.99) (48,-0.76) (49,-0.26) (50,-0.58) (51,-0.67) (52,0.36) (53,1.45) (54,-0.92) (55,0.53) (56,-0.79) (57,1.02) (58,-1.25) (59,-2.03) (60,-1.65) (61,0.38) (62,-0.23) (63,0.68) (64,0.07) (65,0.22) (66,0.14) (67,0.84) (68,-1.56) (69,-0.46) (70,1.37) (71,-0.00) (72,1.52) (73,0.38) (74,0.09) (75,0.17) (76,-0.61) (77,-1.01) (78,1.84) (79,1.17) (80,1.96) (81,1.25) (82,0.39) (83,0.04) (84,0.39) (85,0.15) (86,0.33) (87,-0.45) (88,-0.22) (89,0.06) (90,1.60) (91,1.40) (92,1.06) (93,-0.24) (94,-0.28) (95,1.11) (96,0.58) (97,-0.43) (98,-0.03) (99,1.10) (100,-0.61) (101,1.17) (102,1.13) (103,1.44) (104,-1.08) (105,0.67) (106,0.96) (107,0.16) (108,-0.28) (109,0.73) (110,-0.61) (111,0.02) (112,0.31) (113,1.12) (114,-1.21) (115,1.28) (116,1.11) (117,-1.02) (118,-0.65) (119,-0.11) (120,1.35) (121,-0.37) (122,0.05) (123,1.07) (124,1.23) (125,0.18) (126,1.19) (127,-1.00) (128,1.50) (129,-2.04) (130,0.19) (131,0.12) (132,-0.27) (133,-0.83) (134,-0.72) (135,0.07) (136,0.31) (137,-1.17) (138,-0.87) (139,0.80) (140,-0.65) (141,-0.58) (142,-0.77) (143,0.11) (144,0.13) (145,-0.35) (146,0.45) (147,-0.64) (148,-0.42) (149,0.03) (150,-0.49) (151,0.07) (152,-0.06) (153,0.71) (154,-0.87) (155,1.05) (156,-0.33) (157,0.83) (158,0.99) (159,1.23) (160,-0.62) (161,-0.47) (162,0.55) (163,1.57) (164,1.27) (165,1.43) (166,-0.07) (167,1.41) (168,1.25) (169,0.50) (170,0.89) (171,2.44) (172,1.58) (173,-0.88) (174,-0.48) (175,-0.39) (176,0.41) (177,-1.21) (178,2.38) (179,2.61) (180,-2.52) (181,1.51) (182,-3.25) (183,5.18) (184,6.14) (185,1.32) (186,-0.14) (187,0.17) (188,1.25) (189,-0.53) (190,-1.29) (191,-2.67) (192,1.47) (193,0.89) (194,0.43) (195,-0.94) (196,-0.76) (197,-0.01) (198,-1.57) (199,0.66) (200,0.08) (201,0.26) (202,-0.65) (203,0.16) (204,-0.55) (205,-1.76) (206,1.02) (207,0.39) (208,0.04) (209,1.14) (210,0.45) (211,-0.91) (212,0.43) (213,-0.27) (214,-0.60) (215,1.74) (216,0.32) (217,0.36) (218,0.85) (219,0.93) (220,2.08) (221,0.20) (222,-0.29) (223,-0.53) (224,0.37) (225,0.03) (226,0.22) (227,0.43) (228,0.54) (229,-0.02) (230,-0.17) (231,1.02) (232,1.29) (233,-0.48) (234,0.80) (235,0.96) (236,-0.25) (237,1.25) (238,-0.54) (239,-1.01) (240,-2.09) (241,1.37) (242,1.08) (243,-1.01) (244,0.82) (245,-1.16) (246,-1.33) (247,-0.78) (248,-0.84) (249,-1.46) (250,1.01) (251,0.23) (252,0.28) (253,0.74) (254,-1.35) (255,1.51) (256,0.33) (257,-0.69) (258,0.09) (259,-0.50) (260,-0.52) (261,0.38) (262,-0.54) (263,-0.03) (264,0.37) (265,-0.18) (266,0.26) (267,-0.82) (268,-0.05) (269,0.90) (270,1.81) (271,0.34) (272,-0.72) (273,-1.15) (274,-0.78) (275,-1.06) (276,0.14) (277,0.20) (278,-0.21) (279,0.49) (280,-0.06) (281,0.09) (282,-0.52) (283,-1.01) (284,0.03) (285,0.09) (286,0.32) (287,0.72) (288,-0.46) (289,0.14) (290,1.38) (291,-2.92) (292,-1.74) (293,-0.33) (294,0.34) (295,-0.77) (296,0.07) (297,-0.83) (298,1.16) (299,0.32) (300,-1.26) (301,-1.10) (302,0.75) (303,-1.53) (304,-0.09) (305,-0.17) (306,-0.34) (307,0.06) (308,0.85) (309,1.04) (310,0.35) (311,-0.06) (312,-0.07) (313,-0.46) (314,-0.14) (315,-1.15) (316,-0.89) (317,0.09) (318,0.54) (319,0.01) (320,0.57) (321,0.91) (322,0.11) (323,-0.72) (324,0.84) (325,-0.12) (326,0.11) (327,-0.98) (328,-0.12) (329,0.02) (330,0.75) (331,-0.83) (332,1.06) (333,-1.39) (334,0.92) (335,-0.01) (336,-0.48) (337,0.64) (338,-0.04) (339,-0.18) (340,0.28) (341,-0.63) (342,0.78) (343,-1.28) (344,-0.62) (345,0.24) (346,-0.57) (347,0.25) (348,1.22) (349,1.83) (350,-2.48) (351,-0.06) (352,-1.20) (353,-0.59) (354,-0.31) (355,0.31) (356,-0.27) (357,0.93) (358,-1.84) (359,1.45) (360,-2.23) (361,2.12) (362,1.68) (363,0.91) (364,-0.07) (365,-0.82) (366,0.92) (367,0.41) (368,1.02) (369,-0.09) (370,-0.03) (371,0.40) (372,-0.38) (373,0.86) (374,0.98) (375,1.48) (376,0.86) (377,0.97) (378,0.06) (379,-1.27) (380,0.83) (381,-0.04) (382,-0.00) (383,-0.78) (384,0.74) (385,-0.35) (386,0.19) (387,0.25) (388,0.48) (389,0.11) (390,-0.11) (391,-0.24) (392,0.10) (393,-0.10) (394,-0.65) (395,2.11) (396,3.90) (397,3.70) (398,1.95) (399,2.18) (400,1.36) (401,-0.21) (402,-0.06) (403,0.12) (404,0.04) (405,-0.30) (406,0.15) (407,0.08) (408,-0.88) (409,-0.74) (410,1.06) (411,0.83) (412,0.47) (413,-0.18) (414,-0.12) (415,-1.16) (416,-0.66) (417,-0.63) (418,-0.73) (419,-0.80) (420,0.31) (421,0.14) (422,-1.00) (423,-0.05) (424,-0.02) (425,0.15) (426,0.30) (427,-0.28) (428,-0.65) (429,0.39) (430,-0.14) (431,-1.45) (432,1.06) (433,0.28) (434,0.28) (435,0.25) (436,0.04) (437,0.43) (438,1.42) (439,0.05) (440,-0.32) (441,-0.38) (442,0.57) (443,-1.08) (444,1.17) (445,2.48) (446,-1.49) (447,-0.12) (448,-1.29) (449,2.02) (450,-1.51) (451,1.44) (452,-0.05) (453,-1.30) (454,-1.76) (455,-1.78) (456,-1.37) (457,0.27) (458,1.28) (459,1.18) (460,2.07) (461,-0.79) (462,0.97) (463,-0.42) (464,0.44) (465,-0.76) (466,0.49) (467,-1.00) (468,0.85) (469,0.80) (470,-1.21) (471,0.35) (472,1.47) (473,0.70) (474,1.03) (475,-0.72) (476,0.10) (477,-0.17) (478,-0.09) (479,-0.15) (480,0.38) (481,-0.01) (482,-0.05) (483,1.11) (484,-1.22) (485,0.65) (486,-0.92) (487,0.57) (488,0.21) (489,-0.02) (490,-0.62) (491,-1.20) (492,-1.16) (493,-0.78) (494,-0.17) (495,-0.62) (496,0.51) (497,-0.49) (498,0.68) (499,0.32) (500,-0.98) (501,-1.47) (502,-0.36) (503,1.29) (504,-1.87) (505,-1.11) (506,-0.29) (507,0.38) (508,-0.27) (509,-0.92) (510,-1.57) (511,-0.25) (512,-0.39) (513,0.35) (514,-0.88) (515,0.13) (516,1.42) (517,-1.56) (518,0.43) (519,0.14) (520,-0.70) (521,-0.20) (522,-0.91) (523,0.01) (524,-1.60) (525,-0.14) (526,-0.76) (527,0.06) (528,1.82) (529,0.31) (530,0.90) (531,-1.13) (532,-0.07) (533,-0.51) (534,0.34) (535,0.10) (536,0.38) (537,0.30) (538,0.60) (539,-1.04) (540,1.94) (541,2.14) (542,0.92) (543,-1.53) (544,1.83) (545,-1.26) (546,1.67) (547,0.56) (548,0.75) (549,0.28) (550,-0.53) (551,0.52) (552,0.41) (553,-0.18) (554,0.68) (555,-0.51) (556,0.50) (557,1.08) (558,0.29) (559,0.36) (560,0.85) (561,0.02) (562,-1.25) (563,0.29) (564,1.46) (565,0.40) (566,-1.15) (567,-0.08) (568,0.32) (569,0.26) (570,-0.67) (571,-0.20) (572,0.32) (573,-0.55) (574,-0.56) (575,-0.20) (576,-0.22) (577,-0.67) (578,-0.66) (579,1.14) (580,1.70) (581,2.01) (582,-0.14) (583,0.15) (584,1.06) (585,1.14) (586,0.72) (587,0.79) (588,1.26) (589,-2.53) (590,-0.11) (591,-0.50) (592,-0.45) (593,0.68) (594,-0.34) (595,0.62) (596,0.15) (597,1.64) (598,-0.72) (599,0.49) (600,1.36) (601,-0.01) (602,1.23) (603,-1.61) (604,1.74) (605,-1.31) (606,-0.95) (607,0.26) (608,-0.32) (609,-1.52) (610,1.57) (611,0.19) (612,0.02) (613,-0.89) (614,-0.30) (615,-0.46) (616,1.60) (617,-1.29) (618,-0.51) (619,0.70) (620,-2.65) (621,-0.76) (622,-1.74) (623,-0.38) (624,-0.00) (625,0.93) (626,0.59) (627,-1.13) (628,0.45) (629,-0.78) (630,0.94) (631,-2.61) (632,-0.51) (633,-0.95) (634,-0.15) (635,-0.48) (636,0.63) (637,-1.64) (638,-2.87) (639,-3.55) (640,3.15) (641,0.53) (642,-1.08) (643,-1.37) (644,0.64) (645,0.94) (646,1.70) (647,0.76) (648,-1.21) (649,1.39) (650,-0.33) (651,-1.77) (652,-2.11) (653,-1.72) (654,0.25) (655,-0.22) (656,-0.11) (657,-0.20) (658,-0.00) (659,-0.25) (660,-0.49) (661,-0.15) (662,0.52) (663,-0.84) (664,0.10) (665,0.22) (666,0.84) (667,-2.09) (668,1.95) (669,0.42) (670,0.55) (671,1.06) (672,0.07) (673,0.35) (674,-0.03) (675,-0.37) (676,0.09) (677,-0.66) (678,0.55) (679,-0.48) (680,-0.83) (681,0.18) (682,-0.38) (683,1.25) (684,-0.25) (685,-0.96) (686,1.20) (687,1.34) (688,-1.99) (689,-1.64) (690,1.68) (691,-0.71) (692,0.77) (693,0.15) (694,-0.25) (695,-1.28) (696,-1.22) (697,-0.44) (698,-0.65) (699,-0.26) (700,-1.38) (701,2.09) (702,0.32) (703,-1.47) (704,2.30) (705,-0.32) (706,-0.52) (707,-0.14) (708,0.07) (709,-1.57) (710,-0.43) (711,0.17) (712,-0.50) (713,0.70) (714,0.15) (715,-0.51) (716,-1.59) (717,1.34) (718,0.21) (719,1.42) (720,-1.37) (721,0.29) (722,0.82) (723,0.82) (724,-0.29) (725,0.52) (726,-0.38) (727,-0.04) (728,1.07) (729,-1.01) (730,0.24) (731,0.76) (732,1.78) (733,-2.14) (734,1.75) (735,0.66) (736,0.56) (737,0.12) (738,0.32) (739,-0.49) (740,-0.43) (741,0.30) (742,-0.32) (743,0.27) (744,-0.41) (745,-0.03) (746,-0.56) (747,1.27) (748,0.65) (749,0.58) (750,-1.34) (751,2.43) (752,-1.10) (753,-1.71) (754,0.94) (755,-0.55) (756,-2.11) (757,0.67) (758,-0.70) (759,-0.67) (760,-0.54) (761,0.77) (762,-0.34) (763,-1.30) (764,0.13) (765,0.26) (766,0.94) (767,-0.40) (768,0.64) (769,-0.02) (770,0.32) (771,1.23) (772,-0.65) (773,-0.52) (774,-0.56) (775,0.57) (776,-0.35) (777,1.36) (778,1.28) (779,2.49) (780,1.39) (781,-0.64) (782,-0.60) (783,0.16) (784,-1.26) (785,0.87) (786,-1.21) (787,-0.82) (788,-2.29) (789,-0.91) (790,-1.74) (791,-0.28) (792,0.89) (793,-1.89) (794,0.98) (795,2.23) (796,-2.17) (797,-0.58) (798,-0.58) (799,0.38) (800,0.01) (801,0.27) (802,0.72) (803,1.36) (804,-0.06) (805,1.26) (806,-1.41) (807,1.35) (808,0.65) (809,-0.41) (810,0.29) (811,0.61) (812,1.29) (813,0.05) (814,0.60) (815,-0.41) (816,0.22) (817,-0.12) (818,-0.55) (819,0.57) (820,0.88) (821,-0.34) (822,0.79) (823,0.75) (824,1.38) (825,-2.23) (826,0.04) (827,-0.66) (828,-1.06) (829,-0.72) (830,1.04) (831,-1.69) (832,1.44) (833,-0.65) (834,-0.80) (835,1.03) (836,-1.50) (837,-0.53) (838,0.04) (839,1.44) (840,-0.56) (841,0.25) (842,-0.17) (843,-0.77) (844,0.88) (845,-0.71) (846,-0.85) (847,-0.82) (848,-0.19) (849,-0.00) (850,-1.27) (851,-0.97) (852,0.94) (853,-3.09) (854,1.10) (855,-2.84) (856,-2.22) (857,2.28) (858,0.84) (859,1.02) (860,0.18) (861,1.20) (862,0.47) (863,-0.15) (864,-0.01) (865,0.45) (866,0.50) (867,-0.21) (868,-0.94) (869,-0.18) (870,-0.86) (871,-0.26) (872,0.90) (873,-1.23) (874,-2.60) (875,-1.59) (876,-1.09) (877,-2.10) (878,0.73) (879,0.28) (880,0.37) (881,-0.02) (882,0.61) (883,-1.11) (884,-1.04) (885,-0.84) (886,0.88) (887,-0.42) (888,-0.89) (889,-0.61) (890,0.57) (891,1.16) (892,-0.44) (893,0.57) (894,-0.17) (895,0.35) (896,1.00) (897,-1.11) (898,-1.13) (899,1.71) (900,-1.60) (901,-3.26) (902,-0.91) (903,-1.11) (904,0.03) (905,0.29) (906,-0.83) (907,-1.58) (908,-1.45) (909,0.03) (910,0.10) (911,0.68) (912,0.63) (913,-0.32) (914,0.80) (915,-0.83) (916,0.14) (917,0.50) (918,-0.30) (919,-0.88) (920,-0.01) (921,-0.20) (922,-0.68) (923,-0.04) (924,-0.69) (925,-0.61) (926,-0.08) (927,0.34) (928,-0.64) (929,-0.44) (930,-1.12) (931,-1.12) (932,-0.05) (933,0.82) (934,0.87) (935,1.16) (936,1.49) (937,2.79) (938,1.27) (939,-0.91) (940,0.96) (941,0.21) (942,0.07) (943,0.09) (944,0.35) (945,-1.04) (946,0.17) (947,-0.86) (948,0.67) (949,-0.24) (950,0.34) (951,-0.78) (952,-0.53) (953,-0.63) (954,0.32) (955,-0.34) (956,-1.06) (957,0.08) (958,0.21) (959,-0.04) (960,1.60) (961,0.48) (962,1.41) (963,0.31) (964,0.95) (965,0.69) (966,0.56) (967,-0.06) (968,0.42) (969,-0.02) (970,-0.57) (971,0.42) (972,0.81) (973,0.41) (974,-0.98) (975,-1.17) (976,-0.36) (977,1.02) (978,-1.85) (979,0.98) (980,-0.20) (981,0.89) (982,0.71) (983,-0.73) (984,0.66) (985,1.59) (986,1.27) (987,-0.79) (988,-0.22) (989,0.15) (990,-0.14) (991,-0.31) (992,-0.30) (993,-0.96) (994,0.49) (995,0.17) (996,-0.16) (997,-0.24) (998,-0.08) (999,0.17)};
  \addlegendentry{Daily return}
  \addplot[feorange, thick, no marks] coordinates {(0,2.00) (1,1.97) (2,1.47) (3,1.42) (4,1.78) (5,2.61) (6,1.41) (7,1.48) (8,1.78) (9,1.56) (10,1.60) (11,1.54) (12,1.73) (13,1.44) (14,1.52) (15,1.99) (16,1.46) (17,1.45) (18,1.42) (19,1.66) (20,1.70) (21,1.41) (22,1.42) (23,2.01) (24,1.41) (25,1.52) (26,1.42) (27,1.42) (28,1.49) (29,1.49) (30,1.50) (31,1.54) (32,1.56) (33,2.15) (34,1.41) (35,1.83) (36,1.51) (37,2.12) (38,3.16) (39,1.65) (40,1.51) (41,1.42) (42,2.03) (43,1.45) (44,2.16) (45,1.50) (46,2.71) (47,1.88) (48,1.99) (49,1.77) (50,1.46) (51,1.63) (52,1.70) (53,1.50) (54,2.49) (55,1.92) (56,1.60) (57,1.80) (58,2.02) (59,2.27) (60,3.20) (61,2.73) (62,1.51) (63,1.45) (64,1.71) (65,1.42) (66,1.45) (67,1.43) (68,1.85) (69,2.62) (70,1.56) (71,2.39) (72,1.41) (73,2.57) (74,1.51) (75,1.42) (76,1.43) (77,1.65) (78,2.01) (79,2.96) (80,2.17) (81,3.11) (82,2.26) (83,1.52) (84,1.42) (85,1.52) (86,1.43) (87,1.49) (88,1.55) (89,1.45) (90,1.42) (91,2.66) (92,2.44) (93,2.06) (94,1.46) (95,1.47) (96,2.11) (97,1.63) (98,1.54) (99,1.41) (100,2.10) (101,1.66) (102,2.18) (103,2.14) (104,2.48) (105,2.08) (106,1.70) (107,1.96) (108,1.43) (109,1.47) (110,1.75) (111,1.65) (112,1.41) (113,1.48) (114,2.12) (115,2.22) (116,2.30) (117,2.12) (118,2.02) (119,1.69) (120,1.42) (121,2.38) (122,1.51) (123,1.42) (124,2.08) (125,2.24) (126,1.44) (127,2.20) (128,2.00) (129,2.54) (130,3.21) (131,1.44) (132,1.42) (133,1.46) (134,1.83) (135,1.74) (136,1.42) (137,1.48) (138,2.18) (139,1.88) (140,1.81) (141,1.69) (142,1.63) (143,1.79) (144,1.42) (145,1.43) (146,1.50) (147,1.55) (148,1.68) (149,1.53) (150,1.41) (151,1.57) (152,1.42) (153,1.42) (154,1.74) (155,1.88) (156,2.05) (157,1.49) (158,1.83) (159,1.99) (160,2.24) (161,1.66) (162,1.56) (163,1.61) (164,2.64) (165,2.29) (166,2.47) (167,1.42) (168,2.45) (169,2.27) (170,1.58) (171,1.89) (172,3.73) (173,2.65) (174,1.89) (175,1.57) (176,1.52) (177,1.53) (178,2.21) (179,3.65) (180,3.95) (181,3.83) (182,2.57) (183,4.81) (184,7.47) (185,8.80) (186,2.34) (187,1.43) (188,1.44) (189,2.26) (190,1.60) (191,2.31) (192,4.03) (193,2.51) (194,1.89) (195,1.54) (196,1.94) (197,1.77) (198,1.41) (199,2.63) (200,1.69) (201,1.42) (202,1.46) (203,1.69) (204,1.43) (205,1.62) (206,2.86) (207,2.02) (208,1.52) (209,1.42) (210,2.15) (211,1.55) (212,1.91) (213,1.54) (214,1.47) (215,1.65) (216,2.83) (217,1.48) (218,1.50) (219,1.85) (220,1.93) (221,3.26) (222,1.44) (223,1.47) (224,1.60) (225,1.51) (226,1.41) (227,1.45) (228,1.54) (229,1.61) (230,1.41) (231,1.43) (232,2.02) (233,2.31) (234,1.57) (235,1.81) (236,1.96) (237,1.46) (238,2.26) (239,1.60) (240,2.01) (241,3.27) (242,2.40) (243,2.08) (244,2.01) (245,1.83) (246,2.16) (247,2.35) (248,1.80) (249,1.84) (250,2.50) (251,2.01) (252,1.45) (253,1.47) (254,1.76) (255,2.37) (256,2.56) (257,1.49) (258,1.72) (259,1.42) (260,1.58) (261,1.60) (262,1.51) (263,1.61) (264,1.41) (265,1.51) (266,1.44) (267,1.46) (268,1.83) (269,1.42) (270,1.90) (271,2.92) (272,1.49) (273,1.75) (274,2.16) (275,1.79) (276,2.06) (277,1.43) (278,1.44) (279,1.44) (280,1.58) (281,1.42) (282,1.42) (283,1.59) (284,2.01) (285,1.41) (286,1.42) (287,1.48) (288,1.74) (289,1.56) (290,1.43) (291,2.41) (292,4.36) (293,2.84) (294,1.49) (295,1.49) (296,1.78) (297,1.42) (298,1.84) (299,2.17) (300,1.49) (301,2.27) (302,2.10) (303,1.77) (304,2.58) (305,1.42) (306,1.43) (307,1.50) (308,1.42) (309,1.85) (310,2.04) (311,1.50) (312,1.42) (313,1.42) (314,1.56) (315,1.43) (316,2.15) (317,1.89) (318,1.42) (319,1.61) (320,1.41) (321,1.63) (322,1.91) (323,1.42) (324,1.74) (325,1.85) (326,1.43) (327,1.42) (328,1.98) (329,1.42) (330,1.41) (331,1.77) (332,1.83) (333,2.06) (334,2.42) (335,1.92) (336,1.41) (337,1.57) (338,1.68) (339,1.42) (340,1.44) (341,1.47) (342,1.67) (343,1.79) (344,2.30) (345,1.66) (346,1.45) (347,1.63) (348,1.46) (349,2.23) (350,2.96) (351,3.79) (352,1.42) (353,2.20) (354,1.64) (355,1.48) (356,1.48) (357,1.46) (358,1.93) (359,2.96) (360,2.49) (361,3.46) (362,3.31) (363,2.76) (364,1.91) (365,1.42) (366,1.83) (367,1.92) (368,1.53) (369,2.02) (370,1.42) (371,1.41) (372,1.52) (373,1.51) (374,1.87) (375,1.98) (376,2.52) (377,1.86) (378,1.97) (379,1.42) (380,2.29) (381,1.84) (382,1.42) (383,1.41) (384,1.79) (385,1.76) (386,1.50) (387,1.44) (388,1.46) (389,1.57) (390,1.42) (391,1.42) (392,1.46) (393,1.42) (394,1.42) (395,1.69) (396,3.30) (397,5.70) (398,5.43) (399,3.10) (400,3.39) (401,2.38) (402,1.45) (403,1.42) (404,1.42) (405,1.42) (406,1.48) (407,1.43) (408,1.42) (409,1.88) (410,1.76) (411,2.06) (412,1.84) (413,1.56) (414,1.44) (415,1.42) (416,2.17) (417,1.70) (418,1.67) (419,1.75) (420,1.81) (421,1.48) (422,1.43) (423,2.00) (424,1.42) (425,1.41) (426,1.43) (427,1.48) (428,1.47) (429,1.69) (430,1.52) (431,1.43) (432,2.49) (433,2.06) (434,1.47) (435,1.47) (436,1.46) (437,1.42) (438,1.54) (439,2.45) (440,1.42) (441,1.49) (442,1.51) (443,1.63) (444,2.08) (445,2.18) (446,3.78) (447,2.53) (448,1.42) (449,2.30) (450,3.19) (451,2.57) (452,2.48) (453,1.42) (454,2.32) (455,2.86) (456,2.88) (457,2.39) (458,1.47) (459,2.29) (460,2.18) (461,3.25) (462,1.80) (463,1.97) (464,1.53) (465,1.55) (466,1.78) (467,1.57) (468,2.00) (469,1.86) (470,1.81) (471,2.22) (472,1.50) (473,2.51) (474,1.73) (475,2.03) (476,1.74) (477,1.42) (478,1.43) (479,1.42) (480,1.43) (481,1.51) (482,1.41) (483,1.42) (484,2.11) (485,2.23) (486,1.69) (487,1.92) (488,1.63) (489,1.45) (490,1.41) (491,1.67) (492,2.21) (493,2.16) (494,1.80) (495,1.43) (496,1.67) (497,1.59) (498,1.58) (499,1.71) (500,1.48) (501,1.98) (502,2.51) (503,1.50) (504,2.30) (505,3.00) (506,2.11) (507,1.47) (508,1.51) (509,1.46) (510,1.92) (511,2.63) (512,1.46) (513,1.52) (514,1.50) (515,1.89) (516,1.43) (517,2.45) (518,2.62) (519,1.54) (520,1.43) (521,1.73) (522,1.44) (523,1.91) (524,1.41) (525,2.67) (526,1.43) (527,1.78) (528,1.42) (529,2.93) (530,1.48) (531,1.90) (532,2.13) (533,1.42) (534,1.59) (535,1.49) (536,1.42) (537,1.51) (538,1.48) (539,1.65) (540,2.04) (541,3.09) (542,3.34) (543,1.92) (544,2.58) (545,2.95) (546,2.28) (547,2.76) (548,1.62) (549,1.77) (550,1.47) (551,1.60) (552,1.59) (553,1.53) (554,1.44) (555,1.71) (556,1.58) (557,1.58) (558,2.08) (559,1.47) (560,1.50) (561,1.85) (562,1.41) (563,2.26) (564,1.47) (565,2.51) (566,1.52) (567,2.16) (568,1.42) (569,1.49) (570,1.46) (571,1.70) (572,1.44) (573,1.48) (574,1.62) (575,1.62) (576,1.44) (577,1.45) (578,1.70) (579,1.69) (580,2.14) (581,2.79) (582,3.17) (583,1.43) (584,1.43) (585,2.06) (586,2.14) (587,1.74) (588,1.80) (589,2.28) (590,3.85) (591,1.42) (592,1.58) (593,1.55) (594,1.71) (595,1.50) (596,1.66) (597,1.43) (598,2.72) (599,1.74) (600,1.58) (601,2.38) (602,1.41) (603,2.24) (604,2.68) (605,2.84) (606,2.33) (607,1.95) (608,1.46) (609,1.48) (610,2.57) (611,2.63) (612,1.44) (613,1.41) (614,1.90) (615,1.48) (616,1.56) (617,2.67) (618,2.31) (619,1.59) (620,1.73) (621,4.00) (622,1.77) (623,2.84) (624,1.51) (625,1.41) (626,1.93) (627,1.64) (628,2.13) (629,1.55) (630,1.79) (631,1.94) (632,3.96) (633,1.59) (634,1.95) (635,1.43) (636,1.57) (637,1.67) (638,2.72) (639,4.30) (640,5.22) (641,4.68) (642,1.60) (643,2.08) (644,2.40) (645,1.68) (646,1.94) (647,2.79) (648,1.77) (649,2.22) (650,2.43) (651,1.49) (652,2.88) (653,3.30) (654,2.81) (655,1.46) (656,1.45) (657,1.42) (658,1.44) (659,1.41) (660,1.46) (661,1.58) (662,1.43) (663,1.59) (664,1.85) (665,1.42) (666,1.45) (667,1.85) (668,3.28) (669,3.10) (670,1.53) (671,1.61) (672,2.06) (673,1.42) (674,1.50) (675,1.41) (676,1.51) (677,1.42) (678,1.69) (679,1.62) (680,1.57) (681,1.84) (682,1.44) (683,1.51) (684,2.26) (685,1.46) (686,1.96) (687,2.21) (688,2.36) (689,3.15) (690,2.71) (691,2.76) (692,1.73) (693,1.78) (694,1.43) (695,1.46) (696,2.29) (697,2.23) (698,1.55) (699,1.69) (700,1.46) (701,2.41) (702,3.27) (703,1.49) (704,2.51) (705,3.54) (706,1.49) (707,1.59) (708,1.43) (709,1.42) (710,2.63) (711,1.54) (712,1.43) (713,1.58) (714,1.73) (715,1.43) (716,1.59) (717,2.65) (718,2.36) (719,1.45) (720,2.46) (721,2.39) (722,1.47) (723,1.83) (724,1.83) (725,1.47) (726,1.60) (727,1.51) (728,1.42) (729,2.07) (730,2.01) (731,1.45) (732,1.77) (733,2.88) (734,3.34) (735,2.85) (736,1.70) (737,1.62) (738,1.42) (739,1.48) (740,1.58) (741,1.54) (742,1.48) (743,1.49) (744,1.47) (745,1.53) (746,1.42) (747,1.62) (748,2.29) (749,1.69) (750,1.63) (751,2.36) (752,3.71) (753,2.10) (754,2.80) (755,1.95) (756,1.62) (757,3.30) (758,1.70) (759,1.73) (760,1.70) (761,1.61) (762,1.79) (763,1.50) (764,2.32) (765,1.43) (766,1.46) (767,1.94) (768,1.52) (769,1.68) (770,1.41) (771,1.48) (772,2.24) (773,1.69) (774,1.60) (775,1.62) (776,1.63) (777,1.50) (778,2.39) (779,2.30) (780,3.79) (781,2.42) (782,1.68) (783,1.65) (784,1.43) (785,2.27) (786,1.88) (787,2.21) (788,1.83) (789,3.54) (790,1.91) (791,2.84) (792,1.47) (793,1.89) (794,3.03) (795,1.98) (796,3.46) (797,3.38) (798,1.63) (799,1.63) (800,1.51) (801,1.41) (802,1.47) (803,1.74) (804,2.39) (805,1.42) (806,2.28) (807,2.44) (808,2.38) (809,1.69) (810,1.53) (811,1.47) (812,1.66) (813,2.30) (814,1.42) (815,1.65) (816,1.53) (817,1.45) (818,1.42) (819,1.62) (820,1.63) (821,1.89) (822,1.49) (823,1.80) (824,1.77) (825,2.41) (826,3.45) (827,1.42) (828,1.69) (829,2.06) (830,1.75) (831,2.04) (832,2.78) (833,2.48) (834,1.68) (835,1.82) (836,2.03) (837,2.56) (838,1.60) (839,1.42) (840,2.49) (841,1.62) (842,1.46) (843,1.43) (844,1.78) (845,1.89) (846,1.73) (847,1.86) (848,1.83) (849,1.44) (850,1.41) (851,2.29) (852,1.97) (853,1.94) (854,4.59) (855,2.10) (856,4.26) (857,3.45) (858,3.53) (859,1.84) (860,2.02) (861,1.44) (862,2.21) (863,1.57) (864,1.43) (865,1.41) (866,1.55) (867,1.58) (868,1.44) (869,1.95) (870,1.44) (871,1.86) (872,1.46) (873,1.90) (874,2.25) (875,3.93) (876,2.66) (877,2.09) (878,3.29) (879,1.75) (880,1.47) (881,1.51) (882,1.41) (883,1.66) (884,2.11) (885,2.04) (886,1.85) (887,1.89) (888,1.54) (889,1.90) (890,1.65) (891,1.63) (892,2.17) (893,1.54) (894,1.63) (895,1.43) (896,1.50) (897,2.00) (898,2.12) (899,2.13) (900,2.80) (901,2.67) (902,4.82) (903,1.91) (904,2.11) (905,1.41) (906,1.47) (907,1.84) (908,2.64) (909,2.49) (910,1.41) (911,1.42) (912,1.71) (913,1.67) (914,1.49) (915,1.81) (916,1.84) (917,1.43) (918,1.58) (919,1.48) (920,1.88) (921,1.41) (922,1.44) (923,1.71) (924,1.42) (925,1.72) (926,1.65) (927,1.42) (928,1.49) (929,1.68) (930,1.54) (931,2.13) (932,2.12) (933,1.42) (934,1.83) (935,1.88) (936,2.17) (937,2.54) (938,4.19) (939,2.29) (940,1.91) (941,1.96) (942,1.45) (943,1.42) (944,1.42) (945,1.50) (946,2.04) (947,1.44) (948,1.87) (949,1.70) (950,1.45) (951,1.49) (952,1.79) (953,1.60) (954,1.67) (955,1.48) (956,1.49) (957,2.06) (958,1.42) (959,1.44) (960,1.42) (961,2.67) (962,1.57) (963,2.45) (964,1.48) (965,1.95) (966,1.72) (967,1.62) (968,1.42) (969,1.53) (970,1.41) (971,1.63) (972,1.53) (973,1.82) (974,1.53) (975,1.98) (976,2.17) (977,1.50) (978,2.02) (979,2.98) (980,1.98) (981,1.44) (982,1.89) (983,1.73) (984,1.75) (985,1.70) (986,2.66) (987,2.28) (988,1.80) (989,1.45) (990,1.43) (991,1.43) (992,1.48) (993,1.48) (994,1.96) (995,1.57) (996,1.43) (997,1.43) (998,1.45) (999,1.42)};
  \addlegendentry{$\pm 2\sigma_t$, two standard deviations}
  \addplot[feorange, thick, forget plot, no marks] coordinates {(0,-2.00) (1,-1.97) (2,-1.47) (3,-1.42) (4,-1.78) (5,-2.61) (6,-1.41) (7,-1.48) (8,-1.78) (9,-1.56) (10,-1.60) (11,-1.54) (12,-1.73) (13,-1.44) (14,-1.52) (15,-1.99) (16,-1.46) (17,-1.45) (18,-1.42) (19,-1.66) (20,-1.70) (21,-1.41) (22,-1.42) (23,-2.01) (24,-1.41) (25,-1.52) (26,-1.42) (27,-1.42) (28,-1.49) (29,-1.49) (30,-1.50) (31,-1.54) (32,-1.56) (33,-2.15) (34,-1.41) (35,-1.83) (36,-1.51) (37,-2.12) (38,-3.16) (39,-1.65) (40,-1.51) (41,-1.42) (42,-2.03) (43,-1.45) (44,-2.16) (45,-1.50) (46,-2.71) (47,-1.88) (48,-1.99) (49,-1.77) (50,-1.46) (51,-1.63) (52,-1.70) (53,-1.50) (54,-2.49) (55,-1.92) (56,-1.60) (57,-1.80) (58,-2.02) (59,-2.27) (60,-3.20) (61,-2.73) (62,-1.51) (63,-1.45) (64,-1.71) (65,-1.42) (66,-1.45) (67,-1.43) (68,-1.85) (69,-2.62) (70,-1.56) (71,-2.39) (72,-1.41) (73,-2.57) (74,-1.51) (75,-1.42) (76,-1.43) (77,-1.65) (78,-2.01) (79,-2.96) (80,-2.17) (81,-3.11) (82,-2.26) (83,-1.52) (84,-1.42) (85,-1.52) (86,-1.43) (87,-1.49) (88,-1.55) (89,-1.45) (90,-1.42) (91,-2.66) (92,-2.44) (93,-2.06) (94,-1.46) (95,-1.47) (96,-2.11) (97,-1.63) (98,-1.54) (99,-1.41) (100,-2.10) (101,-1.66) (102,-2.18) (103,-2.14) (104,-2.48) (105,-2.08) (106,-1.70) (107,-1.96) (108,-1.43) (109,-1.47) (110,-1.75) (111,-1.65) (112,-1.41) (113,-1.48) (114,-2.12) (115,-2.22) (116,-2.30) (117,-2.12) (118,-2.02) (119,-1.69) (120,-1.42) (121,-2.38) (122,-1.51) (123,-1.42) (124,-2.08) (125,-2.24) (126,-1.44) (127,-2.20) (128,-2.00) (129,-2.54) (130,-3.21) (131,-1.44) (132,-1.42) (133,-1.46) (134,-1.83) (135,-1.74) (136,-1.42) (137,-1.48) (138,-2.18) (139,-1.88) (140,-1.81) (141,-1.69) (142,-1.63) (143,-1.79) (144,-1.42) (145,-1.43) (146,-1.50) (147,-1.55) (148,-1.68) (149,-1.53) (150,-1.41) (151,-1.57) (152,-1.42) (153,-1.42) (154,-1.74) (155,-1.88) (156,-2.05) (157,-1.49) (158,-1.83) (159,-1.99) (160,-2.24) (161,-1.66) (162,-1.56) (163,-1.61) (164,-2.64) (165,-2.29) (166,-2.47) (167,-1.42) (168,-2.45) (169,-2.27) (170,-1.58) (171,-1.89) (172,-3.73) (173,-2.65) (174,-1.89) (175,-1.57) (176,-1.52) (177,-1.53) (178,-2.21) (179,-3.65) (180,-3.95) (181,-3.83) (182,-2.57) (183,-4.81) (184,-7.47) (185,-8.80) (186,-2.34) (187,-1.43) (188,-1.44) (189,-2.26) (190,-1.60) (191,-2.31) (192,-4.03) (193,-2.51) (194,-1.89) (195,-1.54) (196,-1.94) (197,-1.77) (198,-1.41) (199,-2.63) (200,-1.69) (201,-1.42) (202,-1.46) (203,-1.69) (204,-1.43) (205,-1.62) (206,-2.86) (207,-2.02) (208,-1.52) (209,-1.42) (210,-2.15) (211,-1.55) (212,-1.91) (213,-1.54) (214,-1.47) (215,-1.65) (216,-2.83) (217,-1.48) (218,-1.50) (219,-1.85) (220,-1.93) (221,-3.26) (222,-1.44) (223,-1.47) (224,-1.60) (225,-1.51) (226,-1.41) (227,-1.45) (228,-1.54) (229,-1.61) (230,-1.41) (231,-1.43) (232,-2.02) (233,-2.31) (234,-1.57) (235,-1.81) (236,-1.96) (237,-1.46) (238,-2.26) (239,-1.60) (240,-2.01) (241,-3.27) (242,-2.40) (243,-2.08) (244,-2.01) (245,-1.83) (246,-2.16) (247,-2.35) (248,-1.80) (249,-1.84) (250,-2.50) (251,-2.01) (252,-1.45) (253,-1.47) (254,-1.76) (255,-2.37) (256,-2.56) (257,-1.49) (258,-1.72) (259,-1.42) (260,-1.58) (261,-1.60) (262,-1.51) (263,-1.61) (264,-1.41) (265,-1.51) (266,-1.44) (267,-1.46) (268,-1.83) (269,-1.42) (270,-1.90) (271,-2.92) (272,-1.49) (273,-1.75) (274,-2.16) (275,-1.79) (276,-2.06) (277,-1.43) (278,-1.44) (279,-1.44) (280,-1.58) (281,-1.42) (282,-1.42) (283,-1.59) (284,-2.01) (285,-1.41) (286,-1.42) (287,-1.48) (288,-1.74) (289,-1.56) (290,-1.43) (291,-2.41) (292,-4.36) (293,-2.84) (294,-1.49) (295,-1.49) (296,-1.78) (297,-1.42) (298,-1.84) (299,-2.17) (300,-1.49) (301,-2.27) (302,-2.10) (303,-1.77) (304,-2.58) (305,-1.42) (306,-1.43) (307,-1.50) (308,-1.42) (309,-1.85) (310,-2.04) (311,-1.50) (312,-1.42) (313,-1.42) (314,-1.56) (315,-1.43) (316,-2.15) (317,-1.89) (318,-1.42) (319,-1.61) (320,-1.41) (321,-1.63) (322,-1.91) (323,-1.42) (324,-1.74) (325,-1.85) (326,-1.43) (327,-1.42) (328,-1.98) (329,-1.42) (330,-1.41) (331,-1.77) (332,-1.83) (333,-2.06) (334,-2.42) (335,-1.92) (336,-1.41) (337,-1.57) (338,-1.68) (339,-1.42) (340,-1.44) (341,-1.47) (342,-1.67) (343,-1.79) (344,-2.30) (345,-1.66) (346,-1.45) (347,-1.63) (348,-1.46) (349,-2.23) (350,-2.96) (351,-3.79) (352,-1.42) (353,-2.20) (354,-1.64) (355,-1.48) (356,-1.48) (357,-1.46) (358,-1.93) (359,-2.96) (360,-2.49) (361,-3.46) (362,-3.31) (363,-2.76) (364,-1.91) (365,-1.42) (366,-1.83) (367,-1.92) (368,-1.53) (369,-2.02) (370,-1.42) (371,-1.41) (372,-1.52) (373,-1.51) (374,-1.87) (375,-1.98) (376,-2.52) (377,-1.86) (378,-1.97) (379,-1.42) (380,-2.29) (381,-1.84) (382,-1.42) (383,-1.41) (384,-1.79) (385,-1.76) (386,-1.50) (387,-1.44) (388,-1.46) (389,-1.57) (390,-1.42) (391,-1.42) (392,-1.46) (393,-1.42) (394,-1.42) (395,-1.69) (396,-3.30) (397,-5.70) (398,-5.43) (399,-3.10) (400,-3.39) (401,-2.38) (402,-1.45) (403,-1.42) (404,-1.42) (405,-1.42) (406,-1.48) (407,-1.43) (408,-1.42) (409,-1.88) (410,-1.76) (411,-2.06) (412,-1.84) (413,-1.56) (414,-1.44) (415,-1.42) (416,-2.17) (417,-1.70) (418,-1.67) (419,-1.75) (420,-1.81) (421,-1.48) (422,-1.43) (423,-2.00) (424,-1.42) (425,-1.41) (426,-1.43) (427,-1.48) (428,-1.47) (429,-1.69) (430,-1.52) (431,-1.43) (432,-2.49) (433,-2.06) (434,-1.47) (435,-1.47) (436,-1.46) (437,-1.42) (438,-1.54) (439,-2.45) (440,-1.42) (441,-1.49) (442,-1.51) (443,-1.63) (444,-2.08) (445,-2.18) (446,-3.78) (447,-2.53) (448,-1.42) (449,-2.30) (450,-3.19) (451,-2.57) (452,-2.48) (453,-1.42) (454,-2.32) (455,-2.86) (456,-2.88) (457,-2.39) (458,-1.47) (459,-2.29) (460,-2.18) (461,-3.25) (462,-1.80) (463,-1.97) (464,-1.53) (465,-1.55) (466,-1.78) (467,-1.57) (468,-2.00) (469,-1.86) (470,-1.81) (471,-2.22) (472,-1.50) (473,-2.51) (474,-1.73) (475,-2.03) (476,-1.74) (477,-1.42) (478,-1.43) (479,-1.42) (480,-1.43) (481,-1.51) (482,-1.41) (483,-1.42) (484,-2.11) (485,-2.23) (486,-1.69) (487,-1.92) (488,-1.63) (489,-1.45) (490,-1.41) (491,-1.67) (492,-2.21) (493,-2.16) (494,-1.80) (495,-1.43) (496,-1.67) (497,-1.59) (498,-1.58) (499,-1.71) (500,-1.48) (501,-1.98) (502,-2.51) (503,-1.50) (504,-2.30) (505,-3.00) (506,-2.11) (507,-1.47) (508,-1.51) (509,-1.46) (510,-1.92) (511,-2.63) (512,-1.46) (513,-1.52) (514,-1.50) (515,-1.89) (516,-1.43) (517,-2.45) (518,-2.62) (519,-1.54) (520,-1.43) (521,-1.73) (522,-1.44) (523,-1.91) (524,-1.41) (525,-2.67) (526,-1.43) (527,-1.78) (528,-1.42) (529,-2.93) (530,-1.48) (531,-1.90) (532,-2.13) (533,-1.42) (534,-1.59) (535,-1.49) (536,-1.42) (537,-1.51) (538,-1.48) (539,-1.65) (540,-2.04) (541,-3.09) (542,-3.34) (543,-1.92) (544,-2.58) (545,-2.95) (546,-2.28) (547,-2.76) (548,-1.62) (549,-1.77) (550,-1.47) (551,-1.60) (552,-1.59) (553,-1.53) (554,-1.44) (555,-1.71) (556,-1.58) (557,-1.58) (558,-2.08) (559,-1.47) (560,-1.50) (561,-1.85) (562,-1.41) (563,-2.26) (564,-1.47) (565,-2.51) (566,-1.52) (567,-2.16) (568,-1.42) (569,-1.49) (570,-1.46) (571,-1.70) (572,-1.44) (573,-1.48) (574,-1.62) (575,-1.62) (576,-1.44) (577,-1.45) (578,-1.70) (579,-1.69) (580,-2.14) (581,-2.79) (582,-3.17) (583,-1.43) (584,-1.43) (585,-2.06) (586,-2.14) (587,-1.74) (588,-1.80) (589,-2.28) (590,-3.85) (591,-1.42) (592,-1.58) (593,-1.55) (594,-1.71) (595,-1.50) (596,-1.66) (597,-1.43) (598,-2.72) (599,-1.74) (600,-1.58) (601,-2.38) (602,-1.41) (603,-2.24) (604,-2.68) (605,-2.84) (606,-2.33) (607,-1.95) (608,-1.46) (609,-1.48) (610,-2.57) (611,-2.63) (612,-1.44) (613,-1.41) (614,-1.90) (615,-1.48) (616,-1.56) (617,-2.67) (618,-2.31) (619,-1.59) (620,-1.73) (621,-4.00) (622,-1.77) (623,-2.84) (624,-1.51) (625,-1.41) (626,-1.93) (627,-1.64) (628,-2.13) (629,-1.55) (630,-1.79) (631,-1.94) (632,-3.96) (633,-1.59) (634,-1.95) (635,-1.43) (636,-1.57) (637,-1.67) (638,-2.72) (639,-4.30) (640,-5.22) (641,-4.68) (642,-1.60) (643,-2.08) (644,-2.40) (645,-1.68) (646,-1.94) (647,-2.79) (648,-1.77) (649,-2.22) (650,-2.43) (651,-1.49) (652,-2.88) (653,-3.30) (654,-2.81) (655,-1.46) (656,-1.45) (657,-1.42) (658,-1.44) (659,-1.41) (660,-1.46) (661,-1.58) (662,-1.43) (663,-1.59) (664,-1.85) (665,-1.42) (666,-1.45) (667,-1.85) (668,-3.28) (669,-3.10) (670,-1.53) (671,-1.61) (672,-2.06) (673,-1.42) (674,-1.50) (675,-1.41) (676,-1.51) (677,-1.42) (678,-1.69) (679,-1.62) (680,-1.57) (681,-1.84) (682,-1.44) (683,-1.51) (684,-2.26) (685,-1.46) (686,-1.96) (687,-2.21) (688,-2.36) (689,-3.15) (690,-2.71) (691,-2.76) (692,-1.73) (693,-1.78) (694,-1.43) (695,-1.46) (696,-2.29) (697,-2.23) (698,-1.55) (699,-1.69) (700,-1.46) (701,-2.41) (702,-3.27) (703,-1.49) (704,-2.51) (705,-3.54) (706,-1.49) (707,-1.59) (708,-1.43) (709,-1.42) (710,-2.63) (711,-1.54) (712,-1.43) (713,-1.58) (714,-1.73) (715,-1.43) (716,-1.59) (717,-2.65) (718,-2.36) (719,-1.45) (720,-2.46) (721,-2.39) (722,-1.47) (723,-1.83) (724,-1.83) (725,-1.47) (726,-1.60) (727,-1.51) (728,-1.42) (729,-2.07) (730,-2.01) (731,-1.45) (732,-1.77) (733,-2.88) (734,-3.34) (735,-2.85) (736,-1.70) (737,-1.62) (738,-1.42) (739,-1.48) (740,-1.58) (741,-1.54) (742,-1.48) (743,-1.49) (744,-1.47) (745,-1.53) (746,-1.42) (747,-1.62) (748,-2.29) (749,-1.69) (750,-1.63) (751,-2.36) (752,-3.71) (753,-2.10) (754,-2.80) (755,-1.95) (756,-1.62) (757,-3.30) (758,-1.70) (759,-1.73) (760,-1.70) (761,-1.61) (762,-1.79) (763,-1.50) (764,-2.32) (765,-1.43) (766,-1.46) (767,-1.94) (768,-1.52) (769,-1.68) (770,-1.41) (771,-1.48) (772,-2.24) (773,-1.69) (774,-1.60) (775,-1.62) (776,-1.63) (777,-1.50) (778,-2.39) (779,-2.30) (780,-3.79) (781,-2.42) (782,-1.68) (783,-1.65) (784,-1.43) (785,-2.27) (786,-1.88) (787,-2.21) (788,-1.83) (789,-3.54) (790,-1.91) (791,-2.84) (792,-1.47) (793,-1.89) (794,-3.03) (795,-1.98) (796,-3.46) (797,-3.38) (798,-1.63) (799,-1.63) (800,-1.51) (801,-1.41) (802,-1.47) (803,-1.74) (804,-2.39) (805,-1.42) (806,-2.28) (807,-2.44) (808,-2.38) (809,-1.69) (810,-1.53) (811,-1.47) (812,-1.66) (813,-2.30) (814,-1.42) (815,-1.65) (816,-1.53) (817,-1.45) (818,-1.42) (819,-1.62) (820,-1.63) (821,-1.89) (822,-1.49) (823,-1.80) (824,-1.77) (825,-2.41) (826,-3.45) (827,-1.42) (828,-1.69) (829,-2.06) (830,-1.75) (831,-2.04) (832,-2.78) (833,-2.48) (834,-1.68) (835,-1.82) (836,-2.03) (837,-2.56) (838,-1.60) (839,-1.42) (840,-2.49) (841,-1.62) (842,-1.46) (843,-1.43) (844,-1.78) (845,-1.89) (846,-1.73) (847,-1.86) (848,-1.83) (849,-1.44) (850,-1.41) (851,-2.29) (852,-1.97) (853,-1.94) (854,-4.59) (855,-2.10) (856,-4.26) (857,-3.45) (858,-3.53) (859,-1.84) (860,-2.02) (861,-1.44) (862,-2.21) (863,-1.57) (864,-1.43) (865,-1.41) (866,-1.55) (867,-1.58) (868,-1.44) (869,-1.95) (870,-1.44) (871,-1.86) (872,-1.46) (873,-1.90) (874,-2.25) (875,-3.93) (876,-2.66) (877,-2.09) (878,-3.29) (879,-1.75) (880,-1.47) (881,-1.51) (882,-1.41) (883,-1.66) (884,-2.11) (885,-2.04) (886,-1.85) (887,-1.89) (888,-1.54) (889,-1.90) (890,-1.65) (891,-1.63) (892,-2.17) (893,-1.54) (894,-1.63) (895,-1.43) (896,-1.50) (897,-2.00) (898,-2.12) (899,-2.13) (900,-2.80) (901,-2.67) (902,-4.82) (903,-1.91) (904,-2.11) (905,-1.41) (906,-1.47) (907,-1.84) (908,-2.64) (909,-2.49) (910,-1.41) (911,-1.42) (912,-1.71) (913,-1.67) (914,-1.49) (915,-1.81) (916,-1.84) (917,-1.43) (918,-1.58) (919,-1.48) (920,-1.88) (921,-1.41) (922,-1.44) (923,-1.71) (924,-1.42) (925,-1.72) (926,-1.65) (927,-1.42) (928,-1.49) (929,-1.68) (930,-1.54) (931,-2.13) (932,-2.12) (933,-1.42) (934,-1.83) (935,-1.88) (936,-2.17) (937,-2.54) (938,-4.19) (939,-2.29) (940,-1.91) (941,-1.96) (942,-1.45) (943,-1.42) (944,-1.42) (945,-1.50) (946,-2.04) (947,-1.44) (948,-1.87) (949,-1.70) (950,-1.45) (951,-1.49) (952,-1.79) (953,-1.60) (954,-1.67) (955,-1.48) (956,-1.49) (957,-2.06) (958,-1.42) (959,-1.44) (960,-1.42) (961,-2.67) (962,-1.57) (963,-2.45) (964,-1.48) (965,-1.95) (966,-1.72) (967,-1.62) (968,-1.42) (969,-1.53) (970,-1.41) (971,-1.63) (972,-1.53) (973,-1.82) (974,-1.53) (975,-1.98) (976,-2.17) (977,-1.50) (978,-2.02) (979,-2.98) (980,-1.98) (981,-1.44) (982,-1.89) (983,-1.73) (984,-1.75) (985,-1.70) (986,-2.66) (987,-2.28) (988,-1.80) (989,-1.45) (990,-1.43) (991,-1.43) (992,-1.48) (993,-1.48) (994,-1.96) (995,-1.57) (996,-1.43) (997,-1.43) (998,-1.45) (999,-1.42)};
\end{axis}
\end{tikzpicture}%
}
\caption[ARCH(1): volatility rises only briefly]{ARCH(1): volatility rises only briefly. Simulated daily returns from Engle's ARCH(1) model, $\sigma_t^2 = \omega + 0.5\,r_{t-1}^2$. A large day widens the band, but only yesterday's shock matters, so the band snaps back within a day or two and large days come in short bursts.}
\label{fig:10-cluster-arch}
\end{figure}

\begin{figure}[htbp]
\centering
\scalebox{0.9}{%
\begin{tikzpicture}[font=\small]
\begin{axis}[width=0.78\linewidth, height=0.34\linewidth,
  axis lines=left, axis line style={-stealth},
  xmin=0, xmax=1000, ymin=-9, ymax=9,
  xtick={0,200,...,1000}, ytick={-8,-4,0,4,8},
  ylabel={Return (\%)}, clip=false,
  xlabel={Trading day},
  legend style={draw=none, fill=none, font=\footnotesize, at={(0.5,-0.45)}, anchor=north,
                legend columns=-1, /tikz/every even column/.append style={column sep=0.8em}},
  legend cell align=left,]
  \fill[feorange!12] (axis cs:183,-9) rectangle (axis cs:215,9);
  \fill[feorange!12] (axis cs:632,-9) rectangle (axis cs:676,9);
  \fill[fenavy!7]   (axis cs:290,-9) rectangle (axis cs:470,9);
  \addplot[fenavy!80, thin, no marks] coordinates {(0,-0.75) (1,0.22) (2,-0.09) (3,-0.75) (4,1.23) (5,0.01) (6,0.33) (7,-0.75) (8,-0.38) (9,-0.48) (10,0.37) (11,-0.60) (12,0.14) (13,0.34) (14,-0.79) (15,-0.16) (16,-0.20) (17,-0.05) (18,0.48) (19,0.45) (20,0.00) (21,-0.08) (22,0.74) (23,0.01) (24,-0.29) (25,0.04) (26,0.08) (27,0.22) (28,-0.21) (29,0.22) (30,-0.27) (31,0.28) (32,-0.67) (33,0.00) (34,-0.56) (35,0.21) (36,0.74) (37,-1.04) (38,0.25) (39,-0.28) (40,-0.05) (41,0.82) (42,0.14) (43,-0.94) (44,-0.21) (45,-1.36) (46,0.50) (47,-0.79) (48,-0.58) (49,-0.22) (50,-0.56) (51,-0.57) (52,0.29) (53,1.26) (54,-0.57) (55,0.41) (56,-0.71) (57,0.83) (58,-0.93) (59,-1.40) (60,-0.92) (61,0.25) (62,-0.26) (63,0.75) (64,0.06) (65,0.24) (66,0.14) (67,0.79) (68,-1.17) (69,-0.27) (70,1.30) (71,-0.00) (72,1.69) (73,0.28) (74,0.11) (75,0.20) (76,-0.67) (77,-0.96) (78,1.48) (79,0.73) (80,1.64) (81,0.83) (82,0.35) (83,0.04) (84,0.49) (85,0.16) (86,0.37) (87,-0.46) (88,-0.21) (89,0.06) (90,1.48) (91,0.86) (92,0.72) (93,-0.19) (94,-0.30) (95,1.11) (96,0.44) (97,-0.41) (98,-0.02) (99,1.08) (100,-0.44) (101,1.04) (102,0.82) (103,1.07) (104,-0.73) (105,0.54) (106,0.90) (107,0.13) (108,-0.30) (109,0.73) (110,-0.51) (111,0.02) (112,0.30) (113,0.98) (114,-0.81) (115,0.84) (116,0.73) (117,-0.72) (118,-0.49) (119,-0.10) (120,1.32) (121,-0.25) (122,0.05) (123,1.09) (124,0.92) (125,0.13) (126,1.26) (127,-0.77) (128,1.26) (129,-1.45) (130,0.12) (131,0.15) (132,-0.33) (133,-0.93) (134,-0.66) (135,0.07) (136,0.34) (137,-1.17) (138,-0.65) (139,0.68) (140,-0.57) (141,-0.53) (142,-0.70) (143,0.09) (144,0.13) (145,-0.32) (146,0.39) (147,-0.51) (148,-0.31) (149,0.03) (150,-0.39) (151,0.05) (152,-0.04) (153,0.52) (154,-0.53) (155,0.61) (156,-0.18) (157,0.60) (158,0.61) (159,0.72) (160,-0.34) (161,-0.33) (162,0.41) (163,1.11) (164,0.65) (165,0.85) (166,-0.04) (167,1.34) (168,0.82) (169,0.35) (170,0.86) (171,2.03) (172,0.88) (173,-0.68) (174,-0.50) (175,-0.47) (176,0.48) (177,-1.34) (178,2.00) (179,1.61) (180,-1.53) (181,0.98) (182,-3.07) (183,3.40) (184,3.14) (185,0.63) (186,-0.23) (187,0.44) (188,2.91) (189,-0.88) (190,-2.84) (191,-4.46) (192,1.74) (193,1.64) (194,1.00) (195,-2.54) (196,-1.65) (197,-0.03) (198,-4.20) (199,1.17) (200,0.20) (201,0.74) (202,-1.68) (203,0.35) (204,-1.32) (205,-3.61) (206,1.44) (207,0.74) (208,0.09) (209,2.70) (210,0.77) (211,-2.00) (212,0.79) (213,-0.58) (214,-1.26) (215,3.13) (216,0.40) (217,0.80) (218,1.76) (219,1.58) (220,3.38) (221,0.24) (222,-0.71) (223,-1.19) (224,0.72) (225,0.06) (226,0.43) (227,0.76) (228,0.86) (229,-0.03) (230,-0.26) (231,1.46) (232,1.39) (233,-0.48) (234,1.09) (235,1.14) (236,-0.28) (237,1.73) (238,-0.54) (239,-1.35) (240,-2.30) (241,1.11) (242,1.16) (243,-1.24) (244,1.03) (245,-1.55) (246,-1.56) (247,-0.87) (248,-1.17) (249,-1.96) (250,1.09) (251,0.30) (252,0.48) (253,1.16) (254,-1.76) (255,1.58) (256,0.34) (257,-1.11) (258,0.12) (259,-0.78) (260,-0.70) (261,0.49) (262,-0.69) (263,-0.03) (264,0.46) (265,-0.21) (266,0.28) (267,-0.84) (268,-0.04) (269,0.91) (270,1.43) (271,0.20) (272,-0.80) (273,-1.09) (274,-0.63) (275,-1.00) (276,0.11) (277,0.23) (278,-0.22) (279,0.50) (280,-0.05) (281,0.09) (282,-0.46) (283,-0.79) (284,0.02) (285,0.08) (286,0.27) (287,0.55) (288,-0.31) (289,0.10) (290,1.05) (291,-1.55) (292,-0.66) (293,-0.19) (294,0.35) (295,-0.75) (296,0.06) (297,-0.81) (298,0.91) (299,0.22) (300,-1.20) (301,-0.78) (302,0.57) (303,-1.35) (304,-0.06) (305,-0.20) (306,-0.37) (307,0.06) (308,0.84) (309,0.81) (310,0.25) (311,-0.05) (312,-0.07) (313,-0.41) (314,-0.11) (315,-0.94) (316,-0.54) (317,0.06) (318,0.47) (319,0.00) (320,0.46) (321,0.64) (322,0.07) (323,-0.57) (324,0.56) (325,-0.08) (326,0.09) (327,-0.75) (328,-0.07) (329,0.02) (330,0.56) (331,-0.51) (332,0.65) (333,-0.78) (334,0.47) (335,-0.01) (336,-0.40) (337,0.47) (338,-0.03) (339,-0.14) (340,0.21) (341,-0.44) (342,0.48) (343,-0.75) (344,-0.31) (345,0.16) (346,-0.42) (347,0.17) (348,0.87) (349,0.97) (350,-1.11) (351,-0.02) (352,-1.18) (353,-0.42) (354,-0.28) (355,0.30) (356,-0.25) (357,0.83) (358,-1.31) (359,0.79) (360,-1.44) (361,1.12) (362,0.96) (363,0.62) (364,-0.07) (365,-0.99) (366,0.88) (367,0.37) (368,1.12) (369,-0.07) (370,-0.03) (371,0.44) (372,-0.37) (373,0.81) (374,0.77) (375,1.11) (376,0.55) (377,0.81) (378,0.05) (379,-1.34) (380,0.62) (381,-0.03) (382,-0.00) (383,-0.80) (384,0.61) (385,-0.29) (386,0.18) (387,0.23) (388,0.42) (389,0.09) (390,-0.09) (391,-0.19) (392,0.07) (393,-0.07) (394,-0.46) (395,1.27) (396,1.59) (397,1.11) (398,0.64) (399,1.22) (400,0.74) (401,-0.16) (402,-0.07) (403,0.14) (404,0.04) (405,-0.30) (406,0.14) (407,0.07) (408,-0.75) (409,-0.50) (410,0.76) (411,0.53) (412,0.33) (413,-0.14) (414,-0.10) (415,-0.93) (416,-0.39) (417,-0.46) (418,-0.54) (419,-0.56) (420,0.21) (421,0.11) (422,-0.79) (423,-0.03) (424,-0.01) (425,0.12) (426,0.22) (427,-0.20) (428,-0.44) (429,0.23) (430,-0.09) (431,-0.98) (432,0.50) (433,0.16) (434,0.21) (435,0.18) (436,0.03) (437,0.31) (438,0.91) (439,0.02) (440,-0.25) (441,-0.27) (442,0.39) (443,-0.70) (444,0.63) (445,1.33) (446,-0.58) (447,-0.07) (448,-1.23) (449,1.37) (450,-0.84) (451,0.99) (452,-0.03) (453,-1.54) (454,-1.47) (455,-1.30) (456,-1.02) (457,0.24) (458,1.74) (459,1.16) (460,2.13) (461,-0.63) (462,1.32) (463,-0.53) (464,0.67) (465,-1.09) (466,0.61) (467,-1.34) (468,0.93) (469,0.92) (470,-1.40) (471,0.35) (472,2.03) (473,0.67) (474,1.37) (475,-0.84) (476,0.13) (477,-0.25) (478,-0.12) (479,-0.20) (480,0.46) (481,-0.01) (482,-0.05) (483,1.15) (484,-0.93) (485,0.48) (486,-0.86) (487,0.48) (488,0.20) (489,-0.02) (490,-0.61) (491,-0.99) (492,-0.78) (493,-0.55) (494,-0.14) (495,-0.60) (496,0.42) (497,-0.41) (498,0.56) (499,0.24) (500,-0.82) (501,-0.97) (502,-0.20) (503,1.16) (504,-1.24) (505,-0.63) (506,-0.23) (507,0.40) (508,-0.27) (509,-0.89) (510,-1.22) (511,-0.16) (512,-0.42) (513,0.34) (514,-0.84) (515,0.10) (516,1.38) (517,-1.04) (518,0.28) (519,0.15) (520,-0.75) (521,-0.18) (522,-0.91) (523,0.01) (524,-1.61) (525,-0.09) (526,-0.89) (527,0.05) (528,2.03) (529,0.22) (530,1.18) (531,-1.19) (532,-0.07) (533,-0.69) (534,0.39) (535,0.12) (536,0.44) (537,0.31) (538,0.61) (539,-0.94) (540,1.47) (541,1.25) (542,0.53) (543,-1.45) (544,1.42) (545,-0.91) (546,1.53) (547,0.45) (548,0.98) (549,0.33) (550,-0.71) (551,0.62) (552,0.47) (553,-0.20) (554,0.78) (555,-0.48) (556,0.50) (557,1.04) (558,0.22) (559,0.37) (560,0.82) (561,0.02) (562,-1.24) (563,0.20) (564,1.50) (565,0.28) (566,-1.28) (567,-0.07) (568,0.39) (569,0.28) (570,-0.70) (571,-0.18) (572,0.32) (573,-0.51) (574,-0.47) (575,-0.16) (576,-0.19) (577,-0.56) (578,-0.47) (579,0.81) (580,1.02) (581,1.02) (582,-0.07) (583,0.15) (584,1.01) (585,0.82) (586,0.51) (587,0.67) (588,1.03) (589,-1.74) (590,-0.06) (591,-0.64) (592,-0.50) (593,0.73) (594,-0.33) (595,0.65) (596,0.14) (597,1.68) (598,-0.49) (599,0.49) (600,1.44) (601,-0.01) (602,1.53) (603,-1.42) (604,1.37) (605,-1.01) (606,-0.89) (607,0.28) (608,-0.44) (609,-1.93) (610,1.37) (611,0.17) (612,0.02) (613,-1.27) (614,-0.33) (615,-0.61) (616,1.93) (617,-1.08) (618,-0.49) (619,0.93) (620,-3.16) (621,-0.57) (622,-2.75) (623,-0.43) (624,-0.00) (625,1.85) (626,0.89) (627,-1.92) (628,0.62) (629,-1.38) (630,1.44) (631,-3.71) (632,-0.48) (633,-2.06) (634,-0.28) (635,-1.09) (636,1.25) (637,-2.97) (638,-3.73) (639,-3.49) (640,2.81) (641,0.54) (642,-2.97) (643,-3.03) (644,1.26) (645,2.52) (646,3.99) (647,1.38) (648,-3.27) (649,3.16) (650,-0.70) (651,-5.73) (652,-4.44) (653,-3.36) (654,0.57) (655,-0.89) (656,-0.43) (657,-0.70) (658,-0.01) (659,-0.78) (660,-1.39) (661,-0.38) (662,1.31) (663,-1.84) (664,0.19) (665,0.49) (666,1.76) (667,-3.48) (668,2.26) (669,0.52) (670,1.28) (671,2.27) (672,0.13) (673,0.82) (674,-0.05) (675,-0.77) (676,0.16) (677,-1.18) (678,0.83) (679,-0.71) (680,-1.22) (681,0.22) (682,-0.56) (683,1.70) (684,-0.25) (685,-1.40) (686,1.35) (687,1.39) (688,-1.98) (689,-1.35) (690,1.61) (691,-0.69) (692,1.13) (693,0.21) (694,-0.41) (695,-1.91) (696,-1.31) (697,-0.49) (698,-0.98) (699,-0.35) (700,-2.00) (701,2.13) (702,0.27) (703,-2.53) (704,2.73) (705,-0.31) (706,-1.10) (707,-0.26) (708,0.15) (709,-2.90) (710,-0.52) (711,0.32) (712,-0.96) (713,1.19) (714,0.22) (715,-0.87) (716,-2.35) (717,1.40) (718,0.25) (719,2.52) (720,-1.66) (721,0.36) (722,1.56) (723,1.28) (724,-0.44) (725,0.93) (726,-0.60) (727,-0.06) (728,1.66) (729,-1.16) (730,0.28) (731,1.15) (732,2.22) (733,-1.93) (734,1.47) (735,0.65) (736,0.87) (737,0.18) (738,0.53) (739,-0.74) (740,-0.58) (741,0.39) (742,-0.42) (743,0.33) (744,-0.48) (745,-0.04) (746,-0.61) (747,1.19) (748,0.47) (749,0.55) (750,-1.27) (751,1.76) (752,-0.61) (753,-1.60) (754,0.73) (755,-0.59) (756,-2.60) (757,0.53) (758,-1.02) (759,-0.93) (760,-0.74) (761,1.08) (762,-0.43) (763,-1.82) (764,0.13) (765,0.41) (766,1.32) (767,-0.44) (768,0.85) (769,-0.02) (770,0.42) (771,1.45) (772,-0.56) (773,-0.58) (774,-0.63) (775,0.61) (776,-0.37) (777,1.45) (778,0.98) (779,2.00) (780,0.83) (781,-0.58) (782,-0.74) (783,0.19) (784,-1.64) (785,0.81) (786,-1.32) (787,-0.79) (788,-2.60) (789,-0.69) (790,-2.32) (791,-0.28) (792,1.62) (793,-2.76) (794,1.05) (795,3.48) (796,-2.40) (797,-0.67) (798,-1.30) (799,0.81) (800,0.02) (801,0.59) (802,1.40) (803,2.22) (804,-0.07) (805,2.55) (806,-2.00) (807,1.84) (808,0.91) (809,-0.76) (810,0.56) (811,1.17) (812,2.11) (813,0.06) (814,1.17) (815,-0.66) (816,0.36) (817,-0.20) (818,-0.85) (819,0.75) (820,1.12) (821,-0.37) (822,1.03) (823,0.82) (824,1.52) (825,-1.95) (826,0.03) (827,-1.05) (828,-1.40) (829,-0.81) (830,1.33) (831,-1.89) (832,1.30) (833,-0.66) (834,-1.14) (835,1.34) (836,-1.78) (837,-0.53) (838,0.05) (839,2.30) (840,-0.60) (841,0.39) (842,-0.28) (843,-1.18) (844,1.10) (845,-0.83) (846,-1.05) (847,-0.95) (848,-0.22) (849,-0.00) (850,-1.65) (851,-0.88) (852,0.98) (853,-3.23) (854,0.72) (855,-3.83) (856,-1.99) (857,2.52) (858,0.94) (859,2.07) (860,0.34) (861,2.91) (862,0.82) (863,-0.34) (864,-0.01) (865,0.98) (866,0.95) (867,-0.37) (868,-1.72) (869,-0.25) (870,-1.52) (871,-0.36) (872,1.50) (873,-1.63) (874,-3.02) (875,-1.32) (876,-1.30) (877,-3.08) (878,0.80) (879,0.55) (880,0.81) (881,-0.03) (882,1.20) (883,-1.81) (884,-1.39) (885,-1.16) (886,1.31) (887,-0.61) (888,-1.48) (889,-0.83) (890,0.86) (891,1.71) (892,-0.51) (893,0.88) (894,-0.24) (895,0.52) (896,1.37) (897,-1.19) (898,-1.16) (899,1.76) (900,-1.37) (901,-2.98) (902,-0.59) (903,-1.70) (904,0.05) (905,0.56) (906,-1.46) (907,-2.26) (908,-1.61) (909,0.03) (910,0.19) (911,1.21) (912,0.93) (913,-0.47) (914,1.23) (915,-1.06) (916,0.18) (917,0.74) (918,-0.39) (919,-1.14) (920,-0.01) (921,-0.26) (922,-0.82) (923,-0.04) (924,-0.79) (925,-0.57) (926,-0.08) (927,0.35) (928,-0.61) (929,-0.37) (930,-0.99) (931,-0.77) (932,-0.03) (933,0.81) (934,0.69) (935,0.90) (936,1.04) (937,1.77) (938,0.60) (939,-0.75) (940,0.93) (941,0.20) (942,0.09) (943,0.10) (944,0.37) (945,-1.01) (946,0.13) (947,-0.88) (948,0.55) (949,-0.21) (950,0.32) (951,-0.70) (952,-0.41) (953,-0.52) (954,0.25) (955,-0.29) (956,-0.86) (957,0.05) (958,0.18) (959,-0.04) (960,1.27) (961,0.25) (962,1.21) (963,0.20) (964,0.93) (965,0.54) (966,0.48) (967,-0.05) (968,0.40) (969,-0.02) (970,-0.50) (971,0.32) (972,0.64) (973,0.28) (974,-0.76) (975,-0.74) (976,-0.22) (977,0.85) (978,-1.22) (979,0.51) (980,-0.15) (981,0.87) (982,0.55) (983,-0.61) (984,0.54) (985,1.31) (986,0.77) (987,-0.56) (988,-0.19) (989,0.15) (990,-0.14) (991,-0.28) (992,-0.26) (993,-0.79) (994,0.32) (995,0.14) (996,-0.13) (997,-0.19) (998,-0.06) (999,0.13)};
  \addlegendentry{Daily return}
  \addplot[feorange, thick, no marks] coordinates {(0,1.55) (2,1.47) (4,1.41) (6,1.51) (8,1.46) (10,1.36) (12,1.32) (14,1.21) (16,1.23) (18,1.12) (20,1.12) (22,1.03) (24,1.07) (26,1.01) (28,0.96) (30,0.92) (32,0.91) (34,0.98) (36,1.00) (38,1.29) (40,1.19) (42,1.24) (44,1.31) (46,1.54) (48,1.52) (50,1.41) (52,1.37) (54,1.54) (56,1.45) (58,1.50) (60,1.78) (62,1.69) (64,1.59) (66,1.42) (68,1.39) (70,1.48) (72,1.57) (74,1.80) (76,1.59) (78,1.62) (80,1.82) (82,2.02) (84,1.78) (86,1.60) (88,1.47) (90,1.32) (92,1.66) (94,1.55) (96,1.60) (98,1.47) (100,1.52) (102,1.57) (104,1.68) (106,1.60) (108,1.54) (110,1.48) (112,1.36) (114,1.42) (116,1.50) (118,1.51) (120,1.38) (122,1.53) (124,1.56) (126,1.52) (128,1.68) (130,2.00) (132,1.74) (134,1.69) (136,1.55) (138,1.62) (140,1.58) (142,1.49) (144,1.41) (146,1.28) (148,1.24) (150,1.14) (152,1.08) (154,1.06) (156,1.13) (158,1.13) (160,1.23) (162,1.16) (164,1.36) (166,1.43) (168,1.60) (170,1.54) (172,2.08) (174,1.96) (176,1.78) (178,1.86) (180,2.40) (182,2.42) (184,3.81) (186,3.90) (188,3.35) (190,3.53) (192,4.80) (194,4.42) (196,4.25) (198,3.79) (200,4.37) (202,3.77) (204,3.43) (206,4.02) (208,3.61) (210,3.65) (212,3.48) (214,3.06) (216,3.56) (218,3.12) (220,3.14) (222,3.51) (224,3.17) (226,2.76) (228,2.46) (230,2.20) (232,2.19) (234,2.14) (236,2.16) (238,2.26) (240,2.22) (242,2.58) (244,2.52) (246,2.54) (248,2.51) (250,2.70) (252,2.45) (254,2.30) (256,2.57) (258,2.37) (260,2.13) (262,1.94) (264,1.75) (266,1.58) (268,1.53) (270,1.51) (272,1.65) (274,1.74) (276,1.74) (278,1.54) (280,1.42) (282,1.27) (284,1.31) (286,1.18) (288,1.16) (290,1.08) (292,1.65) (294,1.53) (296,1.47) (298,1.44) (300,1.42) (302,1.60) (304,1.76) (306,1.55) (308,1.40) (310,1.48) (312,1.33) (314,1.23) (316,1.31) (318,1.23) (320,1.16) (322,1.19) (324,1.16) (326,1.13) (328,1.16) (330,1.06) (332,1.11) (334,1.24) (336,1.17) (338,1.15) (340,1.05) (342,1.03) (344,1.15) (346,1.08) (348,1.04) (350,1.32) (352,1.40) (354,1.51) (356,1.37) (358,1.37) (360,1.61) (362,1.89) (364,1.83) (366,1.75) (368,1.67) (370,1.65) (372,1.49) (374,1.47) (376,1.61) (378,1.58) (380,1.70) (382,1.55) (384,1.49) (386,1.40) (388,1.27) (390,1.18) (392,1.08) (394,1.00) (396,1.34) (398,1.79) (400,1.85) (402,1.70) (404,1.50) (406,1.35) (408,1.21) (410,1.26) (412,1.30) (414,1.20) (416,1.28) (418,1.23) (420,1.24) (422,1.13) (424,1.16) (426,1.07) (428,1.00) (430,1.00) (432,1.17) (434,1.12) (436,1.05) (438,0.99) (440,1.11) (442,1.05) (444,1.13) (446,1.47) (448,1.36) (450,1.77) (452,1.80) (454,1.93) (456,2.16) (458,2.00) (460,2.26) (462,2.46) (464,2.33) (466,2.21) (468,2.19) (470,2.10) (472,2.07) (474,2.30) (476,2.27) (478,1.98) (480,1.73) (482,1.55) (484,1.60) (486,1.58) (488,1.55) (490,1.38) (492,1.48) (494,1.47) (496,1.38) (498,1.30) (500,1.24) (502,1.43) (504,1.53) (506,1.65) (508,1.49) (510,1.48) (512,1.55) (514,1.42) (516,1.39) (518,1.72) (520,1.53) (522,1.45) (524,1.42) (526,1.66) (528,1.58) (530,1.95) (532,2.06) (534,1.86) (536,1.65) (538,1.50) (540,1.55) (542,1.91) (544,2.00) (546,2.09) (548,2.11) (550,1.96) (552,1.83) (554,1.64) (556,1.57) (558,1.61) (560,1.46) (562,1.41) (564,1.51) (566,1.68) (568,1.71) (570,1.54) (572,1.45) (574,1.36) (576,1.26) (578,1.21) (580,1.28) (582,1.53) (584,1.36) (586,1.51) (588,1.46) (590,1.93) (592,1.75) (594,1.66) (596,1.55) (598,1.83) (600,1.68) (602,1.76) (604,2.11) (606,2.18) (608,1.99) (610,2.24) (612,2.15) (614,2.08) (616,1.87) (618,2.22) (620,2.06) (622,2.80) (624,3.04) (626,2.93) (628,2.94) (630,2.75) (632,3.70) (634,3.51) (636,3.12) (638,3.53) (640,4.65) (642,4.41) (644,4.78) (646,4.55) (648,4.81) (650,5.17) (652,6.06) (654,6.43) (656,5.53) (658,4.75) (660,4.09) (662,3.63) (664,3.49) (666,3.02) (668,3.79) (670,3.60) (672,3.59) (674,3.14) (676,2.75) (678,2.52) (680,2.30) (682,2.16) (684,2.27) (686,2.22) (688,2.35) (690,2.61) (692,2.54) (694,2.33) (696,2.45) (698,2.32) (700,2.13) (702,2.75) (704,2.99) (706,3.16) (708,2.81) (710,3.20) (712,2.78) (714,2.62) (716,2.35) (718,2.76) (720,3.00) (722,2.81) (724,2.79) (726,2.51) (728,2.21) (730,2.36) (732,2.22) (734,2.80) (736,2.64) (738,2.35) (740,2.13) (742,1.91) (744,1.71) (746,1.54) (748,1.66) (750,1.55) (752,2.05) (754,2.16) (756,1.99) (758,2.47) (760,2.34) (762,2.23) (764,2.35) (766,2.06) (768,2.02) (770,1.85) (772,1.95) (774,1.79) (776,1.68) (778,1.83) (780,2.25) (782,2.07) (784,1.87) (786,2.05) (788,2.07) (790,2.54) (792,2.69) (794,3.24) (796,3.81) (798,3.67) (800,3.31) (802,2.87) (804,3.09) (806,3.23) (808,3.35) (810,2.98) (812,2.73) (814,2.74) (816,2.53) (818,2.20) (820,2.06) (822,1.96) (824,1.94) (826,2.42) (828,2.23) (830,2.22) (832,2.52) (834,2.39) (836,2.41) (838,2.42) (840,2.67) (842,2.36) (844,2.22) (846,2.14) (848,2.10) (850,1.83) (852,2.04) (854,3.00) (856,3.81) (858,3.96) (860,3.75) (862,3.85) (864,3.35) (866,2.96) (868,2.63) (870,2.55) (872,2.44) (874,2.61) (876,3.16) (878,3.61) (880,3.16) (882,2.77) (884,2.83) (886,2.74) (888,2.55) (890,2.49) (892,2.54) (894,2.31) (896,2.04) (898,2.17) (900,2.40) (902,3.12) (904,2.98) (906,2.59) (908,2.93) (910,2.74) (912,2.52) (914,2.29) (916,2.27) (918,2.04) (920,1.98) (922,1.73) (924,1.62) (926,1.58) (928,1.42) (930,1.36) (932,1.49) (934,1.44) (936,1.51) (938,1.97) (940,1.85) (942,1.74) (944,1.53) (946,1.56) (948,1.52) (950,1.41) (952,1.37) (954,1.31) (956,1.21) (958,1.24) (960,1.13) (962,1.35) (964,1.46) (966,1.49) (968,1.36) (970,1.25) (972,1.20) (974,1.19) (976,1.31) (978,1.33) (980,1.49) (982,1.47) (984,1.43) (986,1.63) (988,1.58) (990,1.40) (992,1.27) (994,1.29) (996,1.19) (998,1.09) (999,1.05)};
  \addlegendentry{$\pm 2\sigma_t$, two standard deviations}
  \addplot[feorange, thick, no marks, forget plot] coordinates {(0,-1.55) (2,-1.47) (4,-1.41) (6,-1.51) (8,-1.46) (10,-1.36) (12,-1.32) (14,-1.21) (16,-1.23) (18,-1.12) (20,-1.12) (22,-1.03) (24,-1.07) (26,-1.01) (28,-0.96) (30,-0.92) (32,-0.91) (34,-0.98) (36,-1.00) (38,-1.29) (40,-1.19) (42,-1.24) (44,-1.31) (46,-1.54) (48,-1.52) (50,-1.41) (52,-1.37) (54,-1.54) (56,-1.45) (58,-1.50) (60,-1.78) (62,-1.69) (64,-1.59) (66,-1.42) (68,-1.39) (70,-1.48) (72,-1.57) (74,-1.80) (76,-1.59) (78,-1.62) (80,-1.82) (82,-2.02) (84,-1.78) (86,-1.60) (88,-1.47) (90,-1.32) (92,-1.66) (94,-1.55) (96,-1.60) (98,-1.47) (100,-1.52) (102,-1.57) (104,-1.68) (106,-1.60) (108,-1.54) (110,-1.48) (112,-1.36) (114,-1.42) (116,-1.50) (118,-1.51) (120,-1.38) (122,-1.53) (124,-1.56) (126,-1.52) (128,-1.68) (130,-2.00) (132,-1.74) (134,-1.69) (136,-1.55) (138,-1.62) (140,-1.58) (142,-1.49) (144,-1.41) (146,-1.28) (148,-1.24) (150,-1.14) (152,-1.08) (154,-1.06) (156,-1.13) (158,-1.13) (160,-1.23) (162,-1.16) (164,-1.36) (166,-1.43) (168,-1.60) (170,-1.54) (172,-2.08) (174,-1.96) (176,-1.78) (178,-1.86) (180,-2.40) (182,-2.42) (184,-3.81) (186,-3.90) (188,-3.35) (190,-3.53) (192,-4.80) (194,-4.42) (196,-4.25) (198,-3.79) (200,-4.37) (202,-3.77) (204,-3.43) (206,-4.02) (208,-3.61) (210,-3.65) (212,-3.48) (214,-3.06) (216,-3.56) (218,-3.12) (220,-3.14) (222,-3.51) (224,-3.17) (226,-2.76) (228,-2.46) (230,-2.20) (232,-2.19) (234,-2.14) (236,-2.16) (238,-2.26) (240,-2.22) (242,-2.58) (244,-2.52) (246,-2.54) (248,-2.51) (250,-2.70) (252,-2.45) (254,-2.30) (256,-2.57) (258,-2.37) (260,-2.13) (262,-1.94) (264,-1.75) (266,-1.58) (268,-1.53) (270,-1.51) (272,-1.65) (274,-1.74) (276,-1.74) (278,-1.54) (280,-1.42) (282,-1.27) (284,-1.31) (286,-1.18) (288,-1.16) (290,-1.08) (292,-1.65) (294,-1.53) (296,-1.47) (298,-1.44) (300,-1.42) (302,-1.60) (304,-1.76) (306,-1.55) (308,-1.40) (310,-1.48) (312,-1.33) (314,-1.23) (316,-1.31) (318,-1.23) (320,-1.16) (322,-1.19) (324,-1.16) (326,-1.13) (328,-1.16) (330,-1.06) (332,-1.11) (334,-1.24) (336,-1.17) (338,-1.15) (340,-1.05) (342,-1.03) (344,-1.15) (346,-1.08) (348,-1.04) (350,-1.32) (352,-1.40) (354,-1.51) (356,-1.37) (358,-1.37) (360,-1.61) (362,-1.89) (364,-1.83) (366,-1.75) (368,-1.67) (370,-1.65) (372,-1.49) (374,-1.47) (376,-1.61) (378,-1.58) (380,-1.70) (382,-1.55) (384,-1.49) (386,-1.40) (388,-1.27) (390,-1.18) (392,-1.08) (394,-1.00) (396,-1.34) (398,-1.79) (400,-1.85) (402,-1.70) (404,-1.50) (406,-1.35) (408,-1.21) (410,-1.26) (412,-1.30) (414,-1.20) (416,-1.28) (418,-1.23) (420,-1.24) (422,-1.13) (424,-1.16) (426,-1.07) (428,-1.00) (430,-1.00) (432,-1.17) (434,-1.12) (436,-1.05) (438,-0.99) (440,-1.11) (442,-1.05) (444,-1.13) (446,-1.47) (448,-1.36) (450,-1.77) (452,-1.80) (454,-1.93) (456,-2.16) (458,-2.00) (460,-2.26) (462,-2.46) (464,-2.33) (466,-2.21) (468,-2.19) (470,-2.10) (472,-2.07) (474,-2.30) (476,-2.27) (478,-1.98) (480,-1.73) (482,-1.55) (484,-1.60) (486,-1.58) (488,-1.55) (490,-1.38) (492,-1.48) (494,-1.47) (496,-1.38) (498,-1.30) (500,-1.24) (502,-1.43) (504,-1.53) (506,-1.65) (508,-1.49) (510,-1.48) (512,-1.55) (514,-1.42) (516,-1.39) (518,-1.72) (520,-1.53) (522,-1.45) (524,-1.42) (526,-1.66) (528,-1.58) (530,-1.95) (532,-2.06) (534,-1.86) (536,-1.65) (538,-1.50) (540,-1.55) (542,-1.91) (544,-2.00) (546,-2.09) (548,-2.11) (550,-1.96) (552,-1.83) (554,-1.64) (556,-1.57) (558,-1.61) (560,-1.46) (562,-1.41) (564,-1.51) (566,-1.68) (568,-1.71) (570,-1.54) (572,-1.45) (574,-1.36) (576,-1.26) (578,-1.21) (580,-1.28) (582,-1.53) (584,-1.36) (586,-1.51) (588,-1.46) (590,-1.93) (592,-1.75) (594,-1.66) (596,-1.55) (598,-1.83) (600,-1.68) (602,-1.76) (604,-2.11) (606,-2.18) (608,-1.99) (610,-2.24) (612,-2.15) (614,-2.08) (616,-1.87) (618,-2.22) (620,-2.06) (622,-2.80) (624,-3.04) (626,-2.93) (628,-2.94) (630,-2.75) (632,-3.70) (634,-3.51) (636,-3.12) (638,-3.53) (640,-4.65) (642,-4.41) (644,-4.78) (646,-4.55) (648,-4.81) (650,-5.17) (652,-6.06) (654,-6.43) (656,-5.53) (658,-4.75) (660,-4.09) (662,-3.63) (664,-3.49) (666,-3.02) (668,-3.79) (670,-3.60) (672,-3.59) (674,-3.14) (676,-2.75) (678,-2.52) (680,-2.30) (682,-2.16) (684,-2.27) (686,-2.22) (688,-2.35) (690,-2.61) (692,-2.54) (694,-2.33) (696,-2.45) (698,-2.32) (700,-2.13) (702,-2.75) (704,-2.99) (706,-3.16) (708,-2.81) (710,-3.20) (712,-2.78) (714,-2.62) (716,-2.35) (718,-2.76) (720,-3.00) (722,-2.81) (724,-2.79) (726,-2.51) (728,-2.21) (730,-2.36) (732,-2.22) (734,-2.80) (736,-2.64) (738,-2.35) (740,-2.13) (742,-1.91) (744,-1.71) (746,-1.54) (748,-1.66) (750,-1.55) (752,-2.05) (754,-2.16) (756,-1.99) (758,-2.47) (760,-2.34) (762,-2.23) (764,-2.35) (766,-2.06) (768,-2.02) (770,-1.85) (772,-1.95) (774,-1.79) (776,-1.68) (778,-1.83) (780,-2.25) (782,-2.07) (784,-1.87) (786,-2.05) (788,-2.07) (790,-2.54) (792,-2.69) (794,-3.24) (796,-3.81) (798,-3.67) (800,-3.31) (802,-2.87) (804,-3.09) (806,-3.23) (808,-3.35) (810,-2.98) (812,-2.73) (814,-2.74) (816,-2.53) (818,-2.20) (820,-2.06) (822,-1.96) (824,-1.94) (826,-2.42) (828,-2.23) (830,-2.22) (832,-2.52) (834,-2.39) (836,-2.41) (838,-2.42) (840,-2.67) (842,-2.36) (844,-2.22) (846,-2.14) (848,-2.10) (850,-1.83) (852,-2.04) (854,-3.00) (856,-3.81) (858,-3.96) (860,-3.75) (862,-3.85) (864,-3.35) (866,-2.96) (868,-2.63) (870,-2.55) (872,-2.44) (874,-2.61) (876,-3.16) (878,-3.61) (880,-3.16) (882,-2.77) (884,-2.83) (886,-2.74) (888,-2.55) (890,-2.49) (892,-2.54) (894,-2.31) (896,-2.04) (898,-2.17) (900,-2.40) (902,-3.12) (904,-2.98) (906,-2.59) (908,-2.93) (910,-2.74) (912,-2.52) (914,-2.29) (916,-2.27) (918,-2.04) (920,-1.98) (922,-1.73) (924,-1.62) (926,-1.58) (928,-1.42) (930,-1.36) (932,-1.49) (934,-1.44) (936,-1.51) (938,-1.97) (940,-1.85) (942,-1.74) (944,-1.53) (946,-1.56) (948,-1.52) (950,-1.41) (952,-1.37) (954,-1.31) (956,-1.21) (958,-1.24) (960,-1.13) (962,-1.35) (964,-1.46) (966,-1.49) (968,-1.36) (970,-1.25) (972,-1.20) (974,-1.19) (976,-1.31) (978,-1.33) (980,-1.49) (982,-1.47) (984,-1.43) (986,-1.63) (988,-1.58) (990,-1.40) (992,-1.27) (994,-1.29) (996,-1.19) (998,-1.09) (999,-1.05)};
  \node[feorange, font=\footnotesize, anchor=south] at (axis cs:199,9) {turbulent};
  \node[feorange, font=\footnotesize, anchor=south] at (axis cs:654,9) {turbulent};
  \node[fenavy, font=\footnotesize, anchor=south] at (axis cs:380,9) {quiet};
\end{axis}
\end{tikzpicture}%
}
\caption[GARCH(1,1): large days arrive in clusters]{GARCH(1,1): large days arrive in clusters. Simulated daily returns from Bollerslev's GARCH(1,1) model, $\sigma_t^2 = \omega + 0.13\,r_{t-1}^2 + 0.85\,\sigma_{t-1}^2$, with the same shocks and the same average volatility as in Figure~\ref{fig:10-cluster-arch}. Because yesterday's variance carries forward, a large shock keeps the band wide for weeks: large days arrive in turbulent spells (shaded orange) and small days fill long quiet stretches (shaded blue).}
\label{fig:10-cluster-garch}
\end{figure}


The question carries a great deal of weight. Option prices, risk limits, \term{Value at Risk}\index{Value at Risk (VaR)}\label{kt:value-at-risk}, \term{margin} requirements\index{margin} (the deposits that brokers demand from investors who borrow or trade futures) and volatility-targeting\index{volatility targeting} funds all assume that volatility\index{volatility} can be forecast. If it could not, a risk manager would have no better guess for tomorrow's risk than the long-run average, and none of these tools would work as intended.

Yet in the 1990s the evidence looked weak. Researchers judged a GARCH forecast by how well it predicted the next day's squared return, and by that test the forecasts explained only a few percent of the variation. It seemed that the models, though popular, had little real forecasting power.

The question of this chapter is: \textbf{how well can tomorrow's volatility be forecast from today's information, and are the forecasts good enough to improve a portfolio?}

The lab in Section~\ref{sec:10-lab} applies the proxy argument of Andersen and Bollerslev (1998) and the model comparison of Hansen and Lunde (2005) to free daily U.S. market returns since 1926 and the VIX since 1990, and then constructs volatility-managed portfolios in the manner of Moreira and Muir (2017), in order to see how predictable volatility is, which model forecasts it best, and whether the forecasts are useful to an investor.


\section{Foundations}
\label{sec:10-foundations}

\subsection{Volatility}

Volatility, the standard deviation\index{standard deviation} of returns, and its annualization were introduced in Section~\ref{sec:01-expret}. This chapter works mostly with daily data, so a daily standard deviation $\sigma_d$ is annualized as $\sigma_d\sqrt{252}$, since there are about 252 trading days in a year; a 1\% daily volatility is therefore about 16\% a year. 
The square root assumes returns are uncorrelated from day to day, which is close to true for daily stock returns (Section~\ref{sec:10-how}). 

\subsection{Variance}

Variance is volatility squared. 
Most models work with variance rather than volatility, because variances add up over time. If daily returns are uncorrelated, the variance over the next $h$ days is simply the sum of the $h$ daily variances. Volatilities do not add in this way; to get the volatility over $h$ days, you first add the variances and then take the square root.


\subsection{Historical volatility}

\textbf{Historical volatility}\index{volatility!historical|textbf} is the sample standard deviation over a past window, such as the last 22 trading days (about one month). It weights every day in the window equally and drops days abruptly when they leave it. Because daily mean returns are tiny compared with daily volatility, practitioners usually set the mean to zero and estimate variance as the average squared return.

\subsection{EWMA}


\term{EWMA}\index{EWMA|textbf} (exponentially weighted moving average) estimates today's variance as a weighted average of past squared returns. Unlike historical volatility, which gives every day in the window the same weight, EWMA gives the most weight to yesterday's squared return and less to each earlier day, with the weights shrinking by a constant factor $\lambda$ each day. The method was popularized by J.P. Morgan's RiskMetrics\index{RiskMetrics} in 1994.
They used it to estimate the volatilities and correlations that feed into Value-at-Risk (VaR), which tells you how much a portfolio could lose over a short horizon at a given confidence level.
EWMA can be written in a simple recursive form:
\begin{equation}
  \sigma^2_t = \lambda\,\sigma^2_{t-1} + (1 - \lambda)\,r^2_{t-1}
  \label{eq:10-ewma}
\end{equation}
Today's variance estimate $\sigma^2_t$ is a blend of yesterday's estimate $\sigma^2_{t-1}$ and yesterday's squared return $r^2_{t-1}$.
RiskMetrics set $\lambda = 0.94$ for daily data. Yesterday's squared return gets weight 6\%, the day before $0.94 \times 6\% \approx 5.6\%$, and so on, so an old shock fades gradually instead of dropping out all at once (Figure~\ref{fig:10-ewma-shock}).


\begin{figure}[htbp]
\centering
\scalebox{0.9}{%
\begin{tikzpicture}
\begin{axis}[
  width=0.64\linewidth, height=0.36\linewidth,
  axis lines=left, axis line style={-stealth},
  xmin=0, xmax=46, ymin=0.9, ymax=1.75,
  xtick={0,10,...,40}, ytick={1.0,1.2,1.4,1.6},
  yticklabel style={/pgf/number format/fixed,
                    /pgf/number format/precision=1,
                    /pgf/number format/fixed zerofill},
  xlabel={Day},
  ylabel={Estimated daily volatility (\%)},
  clip=false,
  legend style={draw=none, fill=none, at={(0.98,0.98)}, anchor=north east,
                cells={anchor=west}, font=\small},
  legend cell align=left,
]
  % Baseline variance 1 (vol 1%); one 5% move on day 10 adds 25 - 1 = 24
  % EWMA: the shock enters on day 11 with weight 0.06, then decays by 0.94 per day
  \addplot[fenavy, very thick, const plot, domain=0:45, samples=46]
    {sqrt(1 + (x>10.5)*0.06*24*0.94^(x-11))};
  \addlegendentry{EWMA, $\lambda = 0.94$}
  % 22-day window: the shock carries weight 1/22 on days 11--32, then drops out
  \addplot[fegray, very thick, dashed, const plot, domain=0:45, samples=46]
    {sqrt(1 + (x>10.5)*(x<32.5)*24/22)};
  \addlegendentry{22-day window, equal weights}
  % Annotations
  \draw[densely dotted] (axis cs:10,0.9) -- (axis cs:10,1.64)
    node[above, font=\small] {one 5\% move};
  \draw[-stealth, thin] (axis cs:38,1.32)
    node[above, font=\small] {shock leaves the window}
    -- (axis cs:33,1.22);
\end{axis}
\end{tikzpicture}%
}
\caption{EWMA lets an old shock fade gradually instead of dropping it all at once. Response of each estimator to a single 5\% move on day 10 when volatility is otherwise 1\% per day. The 22-day window raises its estimate by a fixed amount for exactly 22 days and then drops back abruptly when the shock leaves the window, even though nothing happened in the market that day. EWMA reacts more strongly at first, because yesterday receives the largest weight, and then lets the shock fade gradually.}
\label{fig:10-ewma-shock}
\end{figure}

\subsection{ARCH}
\label{sec:10-arch}

\textbf{ARCH}\index{ARCH|textbf} (autoregressive conditional heteroskedasticity\index{heteroskedasticity}), introduced by Engle (1982), was the first model in which today's variance depends on recent shocks. In its simplest form, ARCH(1), today's variance depends only on yesterday's squared return:
\begin{equation}
  \sigma^2_t = \omega + \alpha\,r^2_{t-1}
  \label{eq:10-arch}
\end{equation}
Here $\omega > 0$ sets a floor for the variance, and $\alpha$, between 0 and 1, measures how strongly yesterday's shock raises today's variance. A large return yesterday, of either sign, raises today's variance, which makes another large return more likely. This is how the model produces \term{clustering}. The long-run variance is $\omega / (1 - \alpha)$.

ARCH(1) has a short memory. A shock affects only the next day's variance directly, so volatility falls back within a day or two (Figure~\ref{fig:10-cluster-arch}). Real volatility stays high for weeks. To capture this, ARCH($q$) uses the last $q$ squared returns:
\begin{equation}
  \sigma^2_t = \omega + \alpha_1\,r^2_{t-1} + \alpha_2\,r^2_{t-2} + \cdots + \alpha_q\,r^2_{t-q}
  \label{eq:10-archq}
\end{equation}
In practice $q$ had to be large, which meant many parameters to estimate, each of which must be non-negative to keep the variance positive. GARCH solves this problem with a single extra term.


\subsection{GARCH(1,1)}
\label{sec:10-garch}

\textbf{GARCH(1,1)}\index{GARCH|textbf} (generalized autoregressive conditional heteroskedasticity\index{heteroskedasticity}) builds tomorrow's variance from three ingredients:
\begin{equation}
  \sigma^2_t = \omega + \alpha\,r^2_{t-1} + \beta\,\sigma^2_{t-1}
  \label{eq:10-garch}
\end{equation}
Each term has a simple meaning.
\begin{itemize}
  \item $\alpha\,r^2_{t-1}$ is the \emph{news}. A large return yesterday, up or down, raises today's variance. The parameter $\alpha$ sets how strongly the forecast reacts to news.
  \item $\beta\,\sigma^2_{t-1}$ is the \emph{memory}. Part of yesterday's variance carries over to today. The parameter $\beta$ sets how long a shock is remembered.
  \item $\omega$ is a small constant that keeps pulling the variance back toward its normal level.
\end{itemize}
EWMA has news and memory but no $\omega$. GARCH adds this third term, and it is what makes volatility return to normal after a storm.

In plain words, a GARCH forecast jumps after a big day and then drifts back toward a normal level. With typical stock-market values, half of a shock's effect fades in about 34 trading days, or roughly seven weeks. The shaded box shows how these numbers follow from the parameters.

\begin{technote}
\textbf{Technical note: what the GARCH parameters imply.} Four useful quantities follow from the parameters. The examples use $\omega = 0.000002$, $\alpha = 0.08$ and $\beta = 0.90$, the values of the worked example below.
\begin{itemize}
  \item \textbf{Persistence}\index{persistence|textbf}, $\alpha + \beta$, measures how much of today's shock survives to tomorrow. It must be below 1, or volatility would never settle. In the example it is 0.98.
  \item The \term{long-run variance}\index{GARCH!unconditional variance|textbf}, $V = \omega / (1 - \alpha - \beta)$, is the normal level that volatility returns to. In the example $V = 0.0001$, a daily volatility of 1\%, or about 16\% a year.
  \item \textbf{Forecasts revert to $V$ step by step.} Each day ahead, the gap between the forecast and $V$ shrinks by the factor $\alpha + \beta$, so the expected variance $k$ days ahead is
  \begin{equation}
    E_t[\sigma^2_{t+k}] = V + (\alpha + \beta)^{k-1}\left(\sigma^2_{t+1} - V\right)
    \label{eq:10-forecast}
  \end{equation}
  \item The \term{half-life}\index{half-life|textbf} of a shock, the number of days for half of the gap to close, is $\ln 0.5 / \ln(\alpha + \beta)$. In the example it is about 34 trading days (Figure~\ref{fig:10-cluster-garch}).
\end{itemize}

\textbf{EWMA as a special case.} Setting $\omega = 0$, $\alpha = 1 - \lambda$ and $\beta = \lambda$ turns GARCH into EWMA. Its persistence is then exactly 1, so there is no long-run level to return to, and an EWMA forecast stays where it is no matter how far ahead you look.
\end{technote}

\textbf{Falls matter more than rises.} In stock markets, a fall tends to raise volatility more than a rise of the same size. This asymmetry is called the \term{leverage effect}\index{leverage effect|textbf}. \term{GJR-GARCH}\index{GARCH!GJR|textbf} captures it by adding one term to the equation, $\gamma\,r^2_{t-1}\,\mathbf{1}[r_{t-1} < 0]$. The indicator $\mathbf{1}[r_{t-1} < 0]$ equals 1 after a down day and 0 otherwise, so the extra term $\gamma$ applies only after falls.

\textbf{Estimation.} The parameters are estimated by \term{maximum likelihood}\index{maximum likelihood|textbf}, which chooses the values that make the observed returns most probable under the model.


\subsection{Realized variance}
\label{sec:10-rv}

Variance is the quantity we want to forecast, but it has a problem: we never observe it. A day's variance describes how widely the price \emph{could} have moved, yet we see only one outcome, the day's return. The squared daily return is a natural stand-in, but a very noisy one. A risky day can end with the price almost unchanged, and a calm day can end with one large move.

\term{Realized variance}\index{volatility!realized|textbf} solves this by looking inside the day. Instead of one daily return, it uses many short returns, such as the $M$ five-minute returns $r_{t,1}, \dots, r_{t,M}$ within day $t$, and adds up their squares:
\begin{equation}
  RV_t = \sum_{j=1}^{M} r^2_{t,j}
  \label{eq:10-rv}
\end{equation}
Each short return adds a little information about how much the price moved, so the noise averages out. With enough intraday data, realized variance measures the day's variance almost exactly. It turns variance from a hidden quantity into an observed series that we can plot, forecast and use to judge other forecasts. The same idea works at lower frequency: adding up the squared daily returns within a month gives that month's realized variance, which is how the lab builds it from free daily data.

Because realized variance is observed, it can be forecast directly. The \term{HAR model}\index{HAR model|textbf} (heterogeneous autoregressive) regresses tomorrow's realized variance on its average over the last day, week and month:
\begin{equation}
  RV_{t+1} = b_0 + b_d\,RV_t + b_w\,RV^{(w)}_t + b_m\,RV^{(m)}_t + \varepsilon_{t+1}
  \label{eq:10-har}
\end{equation}
where $RV^{(w)}_t$ and $RV^{(m)}_t$ are the averages of daily realized variance over the last 5 and 22 trading days. The three horizons reflect traders who react to news at different speeds, and together they reproduce the long memory of volatility.


\subsection{Implied volatility}
\label{sec:10-iv}

Historical volatility, EWMA, GARCH and HAR are all built from past returns, so they cannot see an event that has not yet moved prices, such as an earnings release, a central-bank meeting or an election next week. \textbf{Option prices can.} An option is worth more when the underlying asset is expected to move more, so option prices contain the market's own view of future volatility.


\term{Implied volatility}\index{volatility!implied|textbf} recovers that view. It is the volatility of the underlying asset that makes an option's model price equal its market price (Section~\ref{sec:14-iv}), so it is the volatility that traders are implicitly paying for today. Because it comes from option prices, it exists only for assets that have traded options: individual stocks, stock indices, major currencies, interest rates, and commodities such as oil and gold. If an asset has no liquid options, there is no implied volatility to recover.

Its use goes beyond option pricing, though. Option traders often quote prices in implied volatility rather than in dollars, because it makes options with different strikes and maturities easier to compare. Outside option trading, people use implied volatility as a forward-looking risk measure. The VIX, for example, is widely watched as a ``fear gauge'' for the stock market. Implied volatility also feeds into volatility forecasts and risk models, and it is traded directly through products such as VIX futures and variance swaps.

\begin{callout}
\textbf{How options work.} An option is a contract that gives you the right, but not the obligation, to buy or sell something at a fixed price before a set date. A \term{call option} gives you the right to buy, and a \textbf{put option} gives you the right to sell. Suppose oil trades at \$80 a barrel and an airline, worried that fuel costs will rise, pays \$4 a barrel for a call that lets it buy oil at \$80 within the next three months. If oil jumps to \$95, the airline can still buy at \$80, which saves \$15 a barrel and leaves it \$11 ahead after the \$4 it paid. If oil falls instead, the airline lets the option expire, buys oil at the cheaper market price and loses only the \$4. A put works the opposite way. An oil producer worried about falling prices can buy a put that lets it sell at \$80, which acts like insurance on its revenue (Figure~\ref{fig:10-option-payoff}). Because the holder gains from large price moves but never loses more than the price paid, an option is worth more when oil is expected to be more volatile. That is why option prices reveal how volatile the market expects oil to be.
\end{callout}


\begin{figure}[htbp]
\centering
\scalebox{0.9}{%
\begin{tikzpicture}
\begin{axis}[
  width=0.64\linewidth, height=0.40\linewidth,
  axis lines=left, axis line style={-stealth},
  xmin=60, xmax=101.5, ymin=-8, ymax=20,
  xtick={60,70,...,100}, ytick={-5,0,5,10,15},
  xlabel={Oil price at expiry (\$ per barrel)},
  ylabel={Profit (\$ per barrel)},
  clip=false,
  legend style={draw=none, fill=none, at={(0.55,1.03)}, anchor=south,
                legend columns=2, font=\small,
                /tikz/every even column/.append style={column sep=1em}},
  legend cell align=left,
]
  % Zero-profit line
  \addplot[black!40, thin, forget plot] coordinates {(60,0) (100,0)};
  % Call: strike 80, premium 4, profit = max(S - 80, 0) - 4
  \addplot[fenavy, very thick] coordinates {(60,-4) (80,-4) (100,16)};
  \addlegendentry{Call (buy at \$80)}
  % Put: strike 80, premium 4, profit = max(80 - S, 0) - 4
  \addplot[fegray, very thick, dashed] coordinates {(60,16) (80,-4) (100,-4)};
  \addlegendentry{Put (sell at \$80)}
  % Strike
  \draw[densely dotted] (axis cs:80,-8) -- (axis cs:80,-4);
  \node[right, font=\small] at (axis cs:80.3,-6.3) {strike \$80};
  % Example outcome for the call
  \fill[fenavy] (axis cs:95,11) circle (2pt);
  \node[left, font=\small, align=right] at (axis cs:94,12.2)
    {oil at \$95:\\ call gains \$11};
  % Maximum loss
  \node[anchor=south west, font=\footnotesize] at (axis cs:60.5,-3.8) {maximum loss: \$4};
\end{axis}
\end{tikzpicture}%
}
\caption{An option's losses are capped and its gains are not, which is why options are worth more when big moves are likely. Profit at expiry per barrel from a call and a put on oil, each with a strike of \$80 and a price of \$4. The call gains when the oil price rises above \$84, the strike plus the price paid, and the put gains when the oil price falls below \$76. If the oil price ends near \$80, either option expires and its holder loses only the \$4. Because the losses are capped and the gains are not, an option is worth more when large price moves in either direction are more likely.}
\label{fig:10-option-payoff}
\end{figure}


The best-known measure of implied volatility is the \term{VIX}\index{VIX|textbf}. It measures the market's expected volatility of the S\&P~500 over the next 30 days, computed from the prices of many index options, and it is quoted in annual percentage points. A VIX of 20 means an expected volatility of about 20\% a year, or about $20/\sqrt{252} \approx 1.3\%$ a day.

Implied volatility has one important bias. In plain words, the VIX tends to run a little high, because it mixes the market's forecast of volatility with the price investors pay for protection. More precisely, it is a forecast under the \emph{\term{risk-neutral}\index{risk-neutral pricing}\label{kt:risk-neutral}} measure, the probabilities that price assets as if investors did not care about risk, not under real-world probabilities. In practice, investors pay extra for options that protect them when volatility jumps, just as they pay for insurance. As a result, the VIX usually exceeds the volatility that actually follows. This gap is the \term{variance risk premium}\index{variance risk premium|textbf}. A forecaster who uses implied volatility must allow for it, and an investor who sells options earns it, which is the subject of Chapter~\ref{ch:19}.

\begin{figure}[htbp]
\centering
\scalebox{0.9}{%
\begin{tikzpicture}
\begin{axis}[
  width=0.88\linewidth, height=0.40\linewidth,
  axis lines=left, axis line style={-stealth},
  xmin=1990, xmax=2026.5, ymin=0, ymax=92,
  xtick={1990,1995,...,2025}, ytick={0,20,40,60,80},
  xticklabel style={/pgf/number format/1000 sep={}},
  ylabel={VIX (annualized \%)},
  clip=false,
]
  % Approximate path built from key levels (not daily data)
  \addplot[fenavy, thick, sharp plot] coordinates {
    (1990.0,20) (1990.6,36) (1991.1,30) (1991.5,17) (1993.0,13) (1994.2,20)
    (1995.0,13) (1996.5,17) (1997.0,19) (1997.8,38) (1998.2,20) (1998.75,45)
    (1999.0,24) (2000.3,30) (2000.9,26) (2001.2,30) (2001.7,43) (2002.1,22)
    (2002.6,45) (2002.8,40) (2003.2,30) (2003.6,20) (2004.5,15) (2005.5,11)
    (2006.4,18) (2006.9,11) (2007.1,11) (2007.6,30) (2008.2,30) (2008.5,22)
    (2008.89,81) (2009.2,50) (2009.6,25) (2010.2,17) (2010.38,46) (2010.9,18)
    (2011.4,16) (2011.6,48) (2011.9,30) (2012.5,17) (2013.5,13) (2014.5,11)
    (2014.8,26) (2015.4,13) (2015.65,41) (2016.1,28) (2016.5,13) (2017.4,11)
    (2017.9,9) (2018.1,37) (2018.6,12) (2018.98,36) (2019.5,15) (2020.0,13)
    (2020.21,83) (2020.5,30) (2020.9,22) (2021.5,17) (2022.2,36) (2022.5,30)
    (2022.8,33) (2023.5,13) (2024.3,14) (2024.59,39) (2024.9,15) (2025.27,52)
    (2025.6,16) (2025.9,17)
  };
  % Long-run average (about 19.5)
  \draw[fegray, dashed] (axis cs:1990,19.5) -- (axis cs:2026.5,19.5)
    node[right, font=\small, text=fegray] {average};
  % Events
  \node[left,  font=\small, align=right] at (axis cs:2008.7,81) {Global financial\\crisis, 2008};
  \node[right, font=\small, align=left]  at (axis cs:2020.4,83) {COVID-19,\\2020};
  \node[above, font=\small] at (axis cs:1998.75,45) {LTCM};
  \node[above, font=\small] at (axis cs:2001.7,43) {9/11};
  \node[above, font=\small] at (axis cs:2011.6,48) {Euro crisis};
  \node[above, font=\small] at (axis cs:2025.27,52) {Tariffs};
\end{axis}
\end{tikzpicture}%
}
\caption{The stock market's ``fear gauge'' stays calm for long stretches and spikes in crises. Approximate path of the VIX, 1990--2025, drawn through its major peaks and troughs rather than daily data. The VIX measures the volatility of the S\&P~500 over the next 30 days implied by index option prices. It spends long periods near or below its average of about 20 and then spikes sharply during crises, the same volatility clustering seen in past returns, now visible in forward-looking prices. Exact daily values are available from FRED (series VIXCLS).}
\label{fig:10-vix}
\end{figure}


\section[Three facts about volatility]{Three facts about volatility that the models must explain}
\label{sec:10-how}

Three regularities in daily stock returns, which Part A of this chapter's lab measures, explain why the models in this chapter take the form they do.

\begin{itemize}
  \item \textbf{Volatility clustering.}\index{volatility clustering|textbf} Returns themselves have almost no autocorrelation\index{autocorrelation}, but absolute and squared returns are positively autocorrelated, and this autocorrelation decays slowly over weeks and months. The slow decay is why ARCH(1) is not enough and why GARCH and HAR, which carry a long memory, fit better. A forecastable size combined with an unforecastable sign is exactly what an efficient market\index{efficient market hypothesis} with changing risk would produce.
  \item \textbf{Fat tails.}\index{fat tails|textbf} Daily returns have far more extreme days than a normal distribution\index{distribution!normal} allows. Part of this comes from clustering itself, because a mixture of calm and turbulent periods, each roughly normal, has fat tails overall. Even after allowing for clustering, the standardized returns $r_t / \sigma_t$ remain somewhat fat-tailed, which is why many applications estimate GARCH with Student-$t$ shocks, drawn from a distribution with fatter tails than the bell curve.
  \item \textbf{The leverage effect.}\index{leverage effect|textbf} Volatility rises more after falls than after rises of the same size, which is why GJR-GARCH adds a term for down days (Section~\ref{sec:10-garch}). Black (1976) attributed this to leverage, since a falling share price raises a firm's debt-to-equity ratio, the size of its borrowing relative to the value of its shares, and so the riskiness of its equity. The effect is too large to be explained by leverage alone, however, and it is strongest for the index. Risk premiums and \term{volatility feedback}\index{volatility feedback|textbf}, in which higher expected volatility lowers prices, also contribute.
\end{itemize}

Lab Figure~\ref{ch:07}.1 shows the first fact on U.S. data. The bars for returns hover near zero at every lag, while the bars for absolute returns are positive at lag 1 and still clearly positive after 50 trading days.


\section{The question}
\label{sec:10-claim}

\begin{callout}
\textbf{The question.} Does today's volatility forecast tomorrow's? Do forecasts built from past returns, or from option prices, beat a constant forecast by a wide margin, out of sample\index{in-sample versus out-of-sample}, at horizons from one day to one month?
\end{callout}

To answer the question, four points need to be settled.

\begin{itemize}
  \item \textbf{What is forecast.} The variance of the market's return over the next $h$ days, for $h = 1$, 5 and 22. True variance is never observed, so we use a \term{proxy}, an observable stand-in: the sum of squared daily returns over those days.
  \item \textbf{Against what.} Beating the \textbf{long-run average}, a constant, shows that volatility is forecastable at all. Beating a \textbf{random walk}\index{random walk|textbf} in volatility, in which last month's variance continues, shows that a model adds something beyond the obvious.
  \item \textbf{How it is scored.} With two \term{loss functions}, MSE\index{loss function!squared-error (MSE)} and QLIKE (Section~\ref{sec:10-eval}).
  \item \textbf{What would count against it.} Forecasts that beat the long-run average only in sample or only in calm periods, a gain that disappears beyond a few days, or no model beating the random walk.
\end{itemize}

A separate question is whether an investor who holds less of the market when forecast volatility is high earns a higher Sharpe ratio\index{Sharpe ratio}. Forecastable volatility is necessary for it but not sufficient, because expected returns\index{return!expected} may rise with volatility and trading has costs.

\section{The original evidence}
\label{sec:10-evidence}

\textbf{ARCH and GARCH.} Engle (1982) introduced the ARCH model, in which today's variance depends on past squared shocks, first applied to U.K. inflation\index{inflation}. Bollerslev (1986) added lagged variance to avoid ARCH's many lags, producing GARCH(1,1). Fitted to almost every financial series, GARCH models found persistence close to 1, so volatility shocks die out over weeks to months. Because stock volatility rises more after falls than after rises, Glosten, Jagannathan and Runkle (1993) added an asymmetric term and found strong leverage effects. Engle shared the 2003 Nobel prize for this work.

\textbf{The low-$R^2$\index{regression!R-squared@$R^2$} puzzle.} Early tests found that GARCH forecasts explained only a few percent of daily squared returns. Andersen and Bollerslev (1998) showed that the fault lay in the test, not the models: a single squared return is a very noisy measure of a day's variance. Measured instead by realized variance\index{realized variance}, the sum of squared returns over short intervals such as five minutes, the same forecasts explained around half of the variation in volatility.

\textbf{Can realized variance itself be forecast?}\index{volatility!realized|textbf} Andersen, Bollerslev, Diebold and Labys (2003) modelled realized variance directly and outperformed GARCH models estimated on daily returns. Corsi (2009) proposed a simpler heterogeneous autoregressive, or HAR, model, which forecasts tomorrow's realized variance from its average over the past day, week and month. It captures volatility's long memory with three regressors and is now the standard benchmark.

\textbf{Which model forecasts best?} Hansen and Lunde (2005) compared 330 models on an exchange rate\index{exchange rate} and IBM stock and found GARCH(1,1) hard to beat, except that models with a leverage effect did better for the stock. Poon and Granger (2003), reviewing 93 studies, found that implied volatility generally forecast best, followed by historical volatility and GARCH, and that all beat a constant.

\textbf{Can investors profit?} Moreira and Muir (2017) noticed that average returns rise much less than volatility, and found that holding less of a portfolio when its recent volatility is high raises its Sharpe ratio. Cederburg, O'Doherty, Wang and Yan (2020) found that across 103 equity strategies the gains mostly disappear when the scaling must be chosen in real time (Section~\ref{sec:10-investors}).

\section{Evaluating volatility forecasts}
\label{sec:10-eval}

Scoring a volatility forecast is harder than scoring a return forecast, because the variance being forecast is never observed, only a noisy proxy. Patton (2011) showed that only some loss functions are \textbf{robust} to that noise: they rank forecasts in the same order as they would against the true variance. Two common ones are robust (Figure~\ref{fig:10-qlike}).

\begin{itemize}
  \item \term{MSE}\index{loss function!squared-error (MSE)|textbf}, $(RV - F)^2$, where $RV$ is the proxy and $F$ the forecast. It is dominated by the largest days, so a few crisis days can decide the comparison.
  \item \term{QLIKE}\index{loss function!QLIKE|textbf}, $\log F + RV/F$. It penalizes forecasts that are too \emph{low} much more than those that are too high, which suits risk management, and it depends on the ratio $RV/F$, so crises do not dominate it.
\end{itemize}

\begin{figure}[htbp]
\centering
\scalebox{0.9}{%
\begin{tikzpicture}[
  % x = forecast / realized variance; both losses measured from their minimum at x = 1
  declare function={qlike(\x) = 1/\x + ln(\x) - 1;  mse(\x) = (\x - 1)^2;}
]
\begin{axis}[
  width=0.61\linewidth, height=0.34\linewidth,
  axis lines=left, axis line style={-stealth},
  xmin=0.2, xmax=3.2, ymin=0, ymax=2.2,
  xtick={0.5,1,...,3}, ytick={0,0.5,1,1.5,2},
  xticklabel style={/pgf/number format/fixed},
  yticklabel style={/pgf/number format/fixed},
  xlabel={Forecast as a multiple of realized variance, $F/RV$},
  ylabel={Loss above its minimum},
  samples=200, clip=false,
]
  \addplot[fegray, thick, dashed, domain=0.2:2.48] {mse(x)};
  \addplot[fenavy, very thick, domain=0.225:3.2] {qlike(x)};
  \node[fegray, anchor=west] at (axis cs:2.5,2.15) {MSE};
  \node[fenavy, anchor=west] at (axis cs:3.22,{qlike(3.2)}) {QLIKE};
  \node[align=center] at (axis cs:0.85,1.4) {too low:\\heavily penalized};
  \node[align=center] at (axis cs:2.65,0.9) {too high:\\mildly penalized};
\end{axis}
\end{tikzpicture}%
}
\caption{For a risk manager, underestimating risk is the costlier mistake, and QLIKE scores forecasts accordingly. Two robust loss functions, measured from their minimum. MSE treats a forecast that is too high and one that is too low symmetrically. QLIKE punishes a forecast that is too low far more than one that is too high: underestimating risk is the costlier mistake.}
\label{fig:10-qlike}
\end{figure}

Mean absolute error, and MSE on volatility rather than variance, are not robust and should be avoided. Two further diagnostics are in the shaded box.

\begin{technote}
\textbf{Technical note: two diagnostics.}
\begin{itemize}
  \item \textbf{Mincer--Zarnowitz regression.} Regress realized on forecast variance. An unbiased forecast has intercept 0 and slope 1. A slope below 1 means the forecast overreacts\index{overreaction}, which is typical of the VIX because of the variance risk premium. The $R^2$ is capped by the noise in the proxy.
  \item \textbf{Diebold--Mariano test}\index{Diebold--Mariano test|textbf} (Chapter~\ref{ch:06}). Tests whether the average loss difference between two models is zero, using Newey--West\index{standard error!Newey--West (HAC)} errors with $h - 1$ lags, since multi-day targets overlap.
\end{itemize}
\end{technote}

Every parameter must be estimated only on data available at the forecast date, as in the walk-forward\index{walk-forward evaluation} design of Chapter~\ref{ch:06}.

\section{Lab~\thechapter: Forecasting Volatility}
\label{sec:10-lab}

\textbf{What you will explore:} how predictable volatility is, which models forecast it best, and whether the forecasts help an investor. Part~A measures the stylized facts of volatility, Part~B follows the proxy argument of Andersen and Bollerslev (1998) and the model comparison of Hansen and Lunde (2005), Part~C repeats that comparison in three crises, and Part~D builds the volatility-managed portfolio of Moreira and Muir (2017). The data are U.S. market returns since 1926 from the daily factors in the Kenneth French\index{Kenneth French data library} library and, from 1990, the VIX from FRED\index{FRED}; both are free.

\textbf{Part A, stylized facts\index{stylized facts of returns}.} Autocorrelations of daily returns and absolute returns at lags 1 to 50, and the correlation\index{correlation} between today's return and tomorrow's absolute return (Lab Figure~\ref{ch:07}.1). Expect the first near zero, the second positive for months, and the third negative: the leverage effect.

\textbf{Part B, forecast comparison.} Each trading day from 1990, seven methods forecast the market's variance over the next 1, 5 and 22 days: the long-run average, a random walk, EWMA with $\lambda = 0.94$, GARCH(1,1), GJR-GARCH, HAR and the squared VIX. Each is scored against the random walk with QLIKE and MSE, Diebold--Mariano tests and Mincer--Zarnowitz regressions (Lab Figures~\ref{ch:07}.2 and~\ref{ch:07}.3). The proxy uses daily returns, because intraday data are not free back to 1926; this makes every model look worse but does not bias the ranking under QLIKE.

\textbf{Part C, crises.} The same losses within 2008--09, 2020 and 2022 (Lab Figure~\ref{ch:07}.4).

\textbf{Part D, use.} A \textbf{volatility-managed portfolio}\index{volatility targeting|textbf} holds the market with weight proportional to 1 / last month's realized variance, as in Moreira and Muir (2017). The lab reports its Sharpe ratio, drawdown\index{drawdown} and alpha\index{alpha}, net of costs and with an optional leverage cap (Lab Figure~\ref{ch:07}.5).

\label{sec:10-stress}
\textbf{Robustness checks.} Every GARCH forecast converges to the long-run level, so beyond a few months a constant is hard to beat. In the first days of October 2008 and March 2020, every forecast built from past returns was far too low, while the VIX, a market price, could react within the day. In 2022 stocks fell by about a fifth while realized volatility stayed moderate, a reminder that volatility and losses are different things. When you have finished, compare what you saw with the published evidence summarized in Section~\ref{sec:10-verdict}.

\section{Where the evidence stands}
\label{sec:10-verdict}

The published studies point to a clear pattern: \textbf{volatility can be forecast well at horizons from one day to one month, but using the forecasts to time the market earns little once they are made in real time.} The table summarizes the evidence as of November 2026.

\begingroup
\small
\setlength{\LTpre}{6pt}\setlength{\LTpost}{6pt}
\begin{longtable}{@{}>{\raggedright\arraybackslash}p{0.19\linewidth}>{\raggedright\arraybackslash}p{0.13\linewidth}>{\raggedright\arraybackslash}p{\dimexpr \linewidth-0.19\linewidth-0.13\linewidth-0.23\linewidth-6\tabcolsep\relax}>{\raggedright\arraybackslash}p{0.23\linewidth}@{}}
\toprule
Situation & Holds? & Why & Evidence \\
\midrule
\endfirsthead
\multicolumn{4}{@{}l}{\emph{(continued)}} \\
\toprule
Situation & Holds? & Why & Evidence \\
\midrule
\endhead
\midrule
\multicolumn{4}{r@{}}{\emph{(continued on the next page)}} \\
\endfoot
\bottomrule
\endlastfoot
Forecasts from one day to one month & \textbf{Yes} & Dynamic models beat the long-run average by a wide margin in real time. & Andersen and Bollerslev (1998); Hansen and Lunde (2005) \\*
\addlinespace[3pt]
Across markets and decades & \textbf{Yes} & Volatility clusters in every liquid market studied, with persistence close to 1. & Bollerslev (1986); Poon and Granger (2003) \\
\addlinespace[3pt]
Judged with the right loss & \textbf{Yes} & Early doubts came from a noisy proxy and an unsuitable loss, both now removed. & Andersen and Bollerslev (1998); Patton (2011) \\
\addlinespace[3pt]
Complex models over simple ones & \textbf{Partly} & GARCH(1,1) is hard to beat; a leverage term helps for stock indices, and HAR helps most with intraday data. & Hansen and Lunde (2005); Andersen, Bollerslev, Diebold and Labys (2003); Corsi (2009) \\
\addlinespace[3pt]
Long horizons & \textbf{No} & Beyond a few months, forecasts converge to an uncertain long-run level. & GARCH mean reversion (Section~\ref{sec:10-garch}) \\
\addlinespace[3pt]
The onset of a crisis & \textbf{No} & Past returns cannot anticipate the first violent day, as in October 2008 and March 2020. & Section~\ref{sec:10-stress} \\
\addlinespace[3pt]
The VIX as a forecast & \textbf{Partly} & It forecasts well but overstates future volatility because of the variance risk premium. & Poon and Granger (2003); Section~\ref{sec:10-iv} \\*
\addlinespace[3pt]
Using forecasts to time the market & \textbf{Weakly} & The gains from volatility-managed portfolios mostly disappear when the scaling is chosen in real time. & Moreira and Muir (2017); Cederburg, O'Doherty, Wang and Yan (2020) \\
\end{longtable}
\addtocounter{table}{-1}% uncaptioned longtable: keep table numbers unchanged
\endgroup

Forecastable volatility predicts risk, not returns, so it does not contradict market efficiency.

\textbf{What would change the picture.} The case for using the forecasts would strengthen if volatility-managed portfolios built only with real-time information kept beating the market after costs and with realistic leverage. The case for forecastability would weaken only if simple models stopped beating the long-run average in new data, which no study so far has found.

\section{What practitioners and investors can use}
\label{sec:10-investors}

\textbf{Risk forecasts, yes.} Anyone sizing positions, setting stop-losses or computing margin should use a forecast that reacts to recent data, such as EWMA or GARCH, not a long-run average. This is the basis of Value at Risk (Chapter~\ref{ch:08}).

\textbf{Volatility timing for returns: be skeptical.} Moreira and Muir (2017) found that scaling market exposure inversely to last month's variance produced an alpha of about 4.9\% a year over 1926--2015, and \textbf{volatility-targeting}\index{volatility targeting|textbf} funds and \textbf{risk-parity}\index{risk parity|textbf} strategies apply versions of the idea. Cederburg, O'Doherty, Wang and Yan (2020), however, found that out of sample the managed versions generally did no better. The managed portfolio also needs leverage\index{leverage} in calm periods, and volatility-targeting funds all sell at once in a spike, which can deepen the declines they try to avoid. The benefit is clearer for momentum\index{momentum}, where crashes follow high volatility (Barroso and Santa-Clara 2015; Chapter~\ref{ch:04}), than for the market.

\begin{callout}
\textbf{What Failed: selling volatility before February 5, 2018.} For years, products that bet against volatility were among the market's best performers. The VelocityShares Daily Inverse VIX Short-Term ETN (XIV) earned the variance risk premium by shorting VIX futures and held \$1.9 billion. On February 5, 2018, the VIX rose 115.6\%, its largest one-day increase ever. XIV lost 96.3\% that day, and Credit Suisse closed it within weeks. Volatility forecasts are informative on average, but a forecast that volatility will stay low is weakest exactly at the onset of a spike, and a strategy that loses almost everything on one day can look excellent in any backtest that misses that day (Chapter~\ref{ch:08}).
\end{callout}

% Key terms of the chapter, in alphabetical order, for readers new to finance.
% A term defined in an earlier chapter points to the chapter, section and page.
\providecolor{FEMaroon}{RGB}{128,30,45}
\providecolor{FECream}{RGB}{253,248,238}
\providecolor{FEGold}{RGB}{196,150,60}
\makeatletter
\@ifundefined{keyterms}{%
  \newtcolorbox{keyterms}[1][Key terms in this chapter]{breakable, enhanced,
    colback=FECream, colframe=FEGold, boxrule=0.6pt, arc=2pt, leftrule=3pt,
    left=6pt, right=6pt, top=3pt, bottom=3pt, before skip=8pt, after skip=8pt,
    fonttitle=\bfseries, coltitle=white, colbacktitle=FEMaroon, title={#1}}%
}{}
\makeatother

\begin{keyterms}
\footnotesize
\begin{description}[leftmargin=0pt, itemsep=0pt, parsep=0pt, topsep=0pt, font=\normalfont\bfseries\color{FEMaroon}]
  \item[ARCH.] ARCH, autoregressive conditional heteroskedasticity, is a model in which today's variance depends on the size of recent returns.
  \item[Call option.] A call option gives its buyer the right to buy an asset at a fixed price before a set date.
  \item[EWMA.] EWMA, the exponentially weighted moving average, estimates today's variance as a weighted average of past squared returns that gives recent days the most weight.
  \item[GARCH.] GARCH, generalized ARCH, lets today's variance depend on both recent squared returns and yesterday's variance, and it is the standard model of financial volatility.
  \item[GJR-GARCH.] GJR-GARCH is a version of GARCH with an extra term that raises the variance more after down days than after up days.
  \item[Half-life.] The half-life of a volatility shock is the number of days it takes for half of the gap between current and long-run variance to close.
  \item[HAR model.] The HAR model forecasts tomorrow's realized variance from its average over the last day, week and month.
  \item[Implied volatility.] Implied volatility is the volatility that makes an option's model price equal its market price, so it reveals the market's own view of future volatility.
  \item[Leverage effect.] The leverage effect is the tendency of volatility to rise more after a fall in prices than after a rise of the same size.
  \item[Long-run variance.] The long-run variance is the normal level of variance to which a GARCH forecast returns over time.
  \item[Loss function.] A loss function scores a forecast by turning its error into a single penalty, so that forecasts can be compared.
  \item[Margin.] Margin is the deposit that a broker demands from an investor who borrows or trades futures.
  \item[Maximum likelihood.] Maximum likelihood estimates a model's parameters by choosing the values that make the observed data most probable.
  \item[MSE.] MSE, the mean squared error, averages the squared forecast errors, so the largest misses dominate it.
  \item[Proxy.] A proxy is an observable stand-in for a quantity that cannot be observed directly, such as summed squared returns for true variance.
  \item[QLIKE.] QLIKE is a loss function for variance forecasts that penalizes forecasts that are too low much more than those that are too high.
  \item[Realized variance.] Realized variance measures a day's variance by adding up the squares of many short intraday returns.
  \item[Risk-neutral measure.] The risk-neutral measure uses adjusted probabilities that price assets as if investors did not care about risk, so every asset is expected to earn the risk-free rate.
  \item[Value at Risk (VaR).] Value at Risk is a loss that should be exceeded only rarely, for example on one trading day in a hundred.
  \item[Variance risk premium.] The variance risk premium is the gap by which implied variance usually exceeds the variance that actually follows, which option sellers earn.
  \item[VIX.] The VIX measures the market's expected volatility of the S\&P~500 over the next 30 days, computed from the prices of index options.
  \item[Volatility clustering.] Volatility clustering is the tendency of large price moves to follow large moves and calm days to follow calm days.
  \item[Volatility feedback.] Volatility feedback is the effect by which higher expected volatility raises the return investors require and so lowers prices.
\end{description}
\end{keyterms}
